\chapter{MPPI 变体与前沿融合：Tube-MPPI、扩散增强、TD-MPC}
\label{ch:mppi-variants}

% ════════════════════════════════════════════════════════════════
%  全吸收章：把 MPPI 基础（\cref{ch:mppi}）之后的三条前沿主线融成一章。
%  吸收 Williams2018(Tube-MPPI) / Gandhi2021(RMPPI) / Kim2022(Smooth-MPPI)
%  + Mohamed2022(Log-MPPI) / Honda2024(SVG-MPPI) / Mohamed2025(U-MPPI)
%  + Yi2024(CoVO-MPC) + Pan2024(MBD) / Xue2025(DIAL-MPC)
%  + Hansen2022(TD-MPC) / Hansen2024(TD-MPC2)。自包含独立记录。
%  本章据论文/讲义重建（个人笔记未同步），数值/论文归属存疑处打 \pz/\rebuilt。
%  教学层遵循 v5.0 理论教学规范。承接前一章 MPPI 基础。
% ════════════════════════════════════════════════════════════════

\begin{practice}[前置自测]
开始之前，先试答下列问题。能流畅作答者可只速读本章表格与速查卡；卡住之处，正是本章要为你解决的。
\begin{enumerate}
  \item \cref{ch:mppi} 的 vanilla MPPI 用一个各向同性高斯撒噪声、开环前向、只看均值。这个朴素设定一次性埋下了哪几类缺陷？
  \item 上真车时仿真里跑得好的 MPPI 常常冲出安全边界。问题出在采样分布的“中心”还是“宽度”？为什么把采样协方差 $\Sigma$ 调大治不了它？
  \item 你听说“扩散模型表达力极强、能生成多峰分布”，但它每一步用的也是高斯。这种表达力从何而来？
  \item 若让 MPPI 只规划很短的时域（如 $3$ 步），却想让它做出“有远见”的决策，缺失的“$3$ 步之后的长期信息”可以从哪里补？
  \item 多条等优的避障路径（向左绕、向右绕）对最优控制分布意味着什么？把它们的控制逐元素平均会发生什么？
\end{enumerate}
\end{practice}

\section*{本章目标}
学完本章，你应当能够：
\begin{enumerate}
  \item \textbf{诊断} vanilla MPPI 的四大缺陷（模型失配脆弱、抖动、杂乱可行性低、单模态）+ 不确定性不可见，并把每个变体一一对应到它要补的那个侧面；
  \item \textbf{说清} Tube-MPPI 的 nominal+ancillary 双层架构与“planner--tracker 冲突”的两个根因，并推出 RMPPI 的两个改动（增广状态、反馈入采样）；
  \item \textbf{陈述} RMPPI 的\emph{自由能增长上界}——MPPI 家族第一个性能保证——及其三个可调部件（约束满足、跟踪性能、采样误差）；
  \item \textbf{掌握} Smooth-MPPI 的 input-lifting（控制导数空间采样）为何\emph{从源头}消除抖动且\emph{不破坏}信息论最优性，及其“扩状态”本质；
  \item \textbf{理解} Log-MPPI 的 NLN 重尾混合分布、SVG-MPPI 的 Stein 变分寻峰（mode-seeking）、U-MPPI 的 Unscented Transform 不确定性传播，并据五侧面表合理选型；
  \item \textbf{证明} MPPI 一步加权更新与 DDPM 一步去噪、与 score 上升三者数学同构，并推出“采样方差 $=$ 平滑尺度、加权平均 $=$ 平滑分布的 score 估计”；
  \item \textbf{说明} MBD（Model-Based Diffusion）如何用已知模型直接算 score、做\emph{免训练免数据}的扩散式轨迹优化，DIAL-MPC 如何用多级退火征服全阶力矩控制；
  \item \textbf{讲清} TD-MPC 把采样规划搬进\emph{学习的 latent 空间}、用\emph{学习的终端价值}让短时域获得长时域远见，以及 TD-MPC2 的规模化（单一超参 + decoder-free 世界模型 + 架构稳定化）。
\end{enumerate}

\subsection*{本章知识导航}
本章是一条“\emph{让采样提议分布越来越接近真实目标分布、再让规划本身越来越强}”的主线。\cref{ch:mppi} 证明过：最优提议分布就是目标（吉布斯）分布 $\mathbb{P}^\star\propto\mathbb{Q}\,e^{-S/\lambda}$，MPPI 及其全部后继做的都是\emph{用各种手段逼近它}。本章三条线（六大变体 / 扩散同构 / 学习世界模型）正是逼近的三个层次。结构如下：

\begin{center}
\small
\begin{tikzpicture}[
  node distance=4.5mm and 7mm, >=Stealth,
  blk/.style={draw, rounded corners, align=center, font=\small, inner sep=3pt, fill=main!6, draw=main!60!black},
  grp/.style={draw, rounded corners, align=center, font=\small, inner sep=3pt, fill=second!6, draw=second!60!black},
  opt/.style={draw, rounded corners, align=center, font=\small, inner sep=3pt, fill=third!8, draw=third!60!black},
  ln/.style={->, thick, gray!60!black}]
\node[blk] (van) {vanilla MPPI\\（朴素高斯提议）};
\node[grp, right=of van] (six) {六大变体\\一侧面一改进};
\node[opt, below=7mm of six] (diff) {扩散同构\\多步去噪 + 退火};
\node[opt, below=7mm of diff] (td) {学习世界模型\\latent + 终端价值};
\node[grp, left=of diff, fill=second!6, draw=second!60!black] (sides) {五侧面\\位置/时间/尾部/模态/不确定};
\draw[ln] (van)--(six);
\draw[ln] (six.south)-- (sides.north);
\draw[ln] (sides)--(diff);
\draw[ln] (six.south) -- (diff.north);
\draw[ln] (diff)--(td);
\end{tikzpicture}
\end{center}

\noindent\textbf{推荐路径}：\cref{sec:mv-six} 是问题驱动的主干——把 vanilla MPPI 的毛病摊成五个侧面，六个变体各补一个（其中 \cref{sec:mv-robust} 的 Tube/RMPPI 理论含量最高、带来家族首个性能保证，值得精读）。\cref{sec:mv-diffusion} 把镜头拉高，揭示“MPPI 一步 $=$ 一步去噪”这个统一同构，并据此得到免训练扩散优化（MBD）与多级退火（DIAL-MPC）。\cref{sec:mv-tdmpc} 走到另一端：不再依赖真实动力学，而是\emph{学}一个 latent 世界模型 + 价值，用熟悉的采样规划内核在其中规划。三节共享同一个采样规划内核，区别只在“提议分布如何被组织”“score 从哪来”“在哪 rollout”。

\subsection*{前置知识桥接}
本章\textbf{承接 \cref{ch:mppi} 的 MPPI 基础}：自由能--KL 推导、重要性采样、零方差最优提议、加权更新公式 $w_k\propto\exp(-(S_k-\rho)/\lambda)$、warm-start 与有效样本数（ESS）、GPU 并行 rollout——这些不再重证，本章直接引用。运动规划的构型空间、采样式与优化式范式见 \cref{ch:planning}；当构型含姿态（$\mathcal{C}=\SE(2)$ 或 $\SE(3)$ 的刚体规划）时，rollout 里的状态插值/距离须用 \cref{ch:lie} 的 $\mathrm{Exp}/\mathrm{Log}$ 与测地线（本章源文献多在 $\mathbb{R}^n$ 状态上叙述，移植到流形状态时按 \cref{ch:lie} 改写）。\cref{sec:mv-tdmpc} 的价值学习用到强化学习的 $Q$ 函数、TD（时序差分）、SAC 双 $Q$ 等基础，正文用到处会用两三句重新激活。

\subsection*{如果跳过本章会怎样}
\begin{itemize}
  \item 你用 \cref{ch:mppi} 的 vanilla MPPI 控越野车，仿真里很好，一上真车摩擦/坡度的建模误差让它冲出安全边界——因为你不知道 \cref{sec:mv-robust}：vanilla 提议建立在“模型准确”的隐含假设上，扰动一来采样中心就失效，盲目调大 $\Sigma$ 治标不治本。
  \item 你读 $2024$ 年后的论文频繁见到“MPPI is a single-step diffusion”“diffusion-style annealing”“latent world model + planning”，却以为这是要训练庞大网络的新领域——不理解 \cref{sec:mv-diffusion,sec:mv-tdmpc}，你会把内核就是 MPPI 的东西当成从头学的陌生方法。
  \item 你想把采样式 MPC 用到几十维力矩控制或像素观测上，单级 MPPI 怎么调都差，只会拼命加样本数撞维度灾难——而退火（\cref{sec:mv-diffusion}）和学习 latent 模型（\cref{sec:mv-tdmpc}）这两把钥匙就在手边。
\end{itemize}

\begin{table}[H]\centering\small
\caption{本章阅读方式建议}
\label{tab:mv-reading}
\begin{tabular}{lll}
\toprule
阅读方式 & 时间 & 适合谁 \\
\midrule
精读（含推导与全部 demo 直觉） & 8--10 小时 & 要实现/选型这些前沿变体的工程师与研究者 \\
速读（抓“缺陷--变体”地图与三大同构） & 3--4 小时 & 有 MPPI 基础、想建立前沿全景 \\
速查（只看五侧面表、选型决策树、速查表） & 30 分钟 & 写代码、选方法时回来查 \\
\bottomrule
\end{tabular}
\end{table}

\begin{note}
\textbf{本章记号与约定.}\;
沿用 \cref{ch:mppi}：控制序列 $U=(u_0,\dots,u_{H-1})$，单条 rollout 噪声 $\varepsilon_k\sim\mathcal{N}(0,\Sigma)$，总代价 $S_k$（或写 $J_k$），温度 $\lambda$，基线 $\rho=\min_k S_k$，权重
\begin{equation}
  w_k=\frac{\exp\!\big(-(S_k-\rho)/\lambda\big)}{\sum_{j=1}^{K}\exp\!\big(-(S_j-\rho)/\lambda\big)},\qquad
  U\leftarrow U+\sum_{k=1}^{K}w_k\,\varepsilon_k.
  \label{eq:mv-mppi-update}
\end{equation}
规划时域 $H$、样本数 $K$、折扣 $\gamma$。\textbf{字母隔离}：本章 $\lambda$ 在 MPPI/扩散语境指\emph{温度}，在 \cref{sec:mv-tdmpc} 的训练损失里临时指\emph{时步衰减权重}（用到处显式声明）；$Q$ 指强化学习\emph{动作价值}（非旋转）；旋转/位姿仍遵全书 $\mathbf{R}\in\SO(3)$、$\mathbf{T}\in\SE(3)$。本章部分数值（误差倍数、命中率）转自源讲义 demo 与原论文，具体以原文为准。
\end{note}

% ════════════════════════════════════════════════════════════════
\section{六大变体：一侧面，一改进}
\label{sec:mv-six}

\paragraph{动机：四大缺陷同出一源。}
\cref{ch:mppi} 交付了一个能部署的 vanilla MPPI，并埋下一条主线——\emph{样本效率}：提议分布是盲目的各向同性高斯，大半 rollout 落在高代价区被浪费。把这条主线往深看一层，会发现 vanilla MPPI“围绕一条控制序列、用一个各向同性高斯撒噪声、开环前向、只看均值”这\emph{一个}设定，同时埋下了四颗雷：高斯是\emph{开环}的（扰动一来就失效$\to$模型失配脆弱）、是\emph{逐点独立}的白噪声（$\to$抖动）、是\emph{薄尾单峰}的（$\to$杂乱环境采不到可行区、多峰取中间）、是\emph{只传均值}的（$\to$看不见不确定性）。

\paragraph{反面：把四缺陷当四个零散 bug 去调。}
若不认清“同根于朴素高斯”，工程上一遇问题就反射性地“调协方差、调温度、加样本”，结果哪个都治不彻底。本节的组织原则因此是\emph{问题驱动}：六个变体不是六个互不相干的改法，而是从同一个朴素高斯出发、各自松开它的\emph{一个}限定——这也是为什么它们大多能正交叠加（松开的是不同限定）、为什么可以用一张“五侧面表”统一（\cref{tab:mv-sides}）。

\begin{insight}
回到 \cref{ch:mppi} 的零方差理想：最优提议分布就是目标分布 $\mathbb{P}^\star\propto\mathbb{Q}\,e^{-S/\lambda}$，而真实环境里这个目标分布常常是\emph{重尾、多峰、形状古怪、闭环、带不确定性}的。单个各向同性高斯从根上拟合不了它。所以六个变体的共同逻辑只有一句：\textbf{目标分布在哪个侧面上最不像高斯，就把提议往那个侧面改一步。} 选型时不必死记六个名字，反过来问“我的 $\mathbb{P}^\star$ 在哪个侧面最不像高斯”即可。
\end{insight}

\subsection{Tube-MPPI 与 RMPPI：闭环鲁棒与首个性能保证}
\label{sec:mv-robust}

\paragraph{动机：扰动把采样中心顶到提议覆盖不到的地方。}
\cref{ch:mppi} 反复强调，MPPI 的成败系于\emph{采样所围绕的那条控制序列}（重要性采样轨迹）；warm-start 有效，是因为上一周期的解是“最优控制长什么样”的好估计。但这套逻辑有一个隐含假设：\emph{系统大致按模型演化}，所以“上一周期的好解”这一周期仍好。真实机器人会打破它——地面摩擦、负载、风、坡度、未建模动力学会把真实状态推到提议\emph{预料之外}的区域\cite{gandhi2021robust}。一旦如此，围绕“过时好解”撒出的 rollout 绝大多数不再描述真实系统，加权平均给出的方向不再可靠。这就是 vanilla MPPI 第二个核心缺陷：\textbf{模型失配脆弱}。

\begin{pitfall}[概念误区：“鲁棒”就是“把协方差调大、多探索”]
遇扰动就把采样协方差 $\Sigma$ 调大，以为撒得更开就能抗扰。\textbf{为何错}：$\Sigma$ 一大，样本效率骤降（\cref{ch:mppi}：大半 rollout 浪费、ESS 崩），且根本没解决“提议\emph{中心}被扰动推偏”的问题——中心错了，撒得再开也是围着错地方撒。这是\emph{位置}问题，不是\emph{宽度}问题。\textbf{正解}：用反馈把真实状态拉回名义附近（让中心重新有效），并让采样在闭环意义下评估，而非盲目加大 $\Sigma$。
\end{pitfall}

\paragraph{反面：让 rollout 看不见扰动会怎样。}
考虑一维系统：状态为位置 $x$ 与速度 $v$，控制 $u$ 为加速度，有一个持续扰动 $d$（如恒定下坡分量）一直推动系统；任务是把 $x$ 停到目标 $4.5$、同时不越墙 $x\le 5$（越过重罚）。若“名义规划”在\emph{无扰动}的想象世界里算出一条恰停在 $4.5$ 的漂亮轨迹，再开环交给带扰动的真实系统，那股一直在推的 $d$ 让速度累积、$x$ 一路冲过墙；而让 rollout \emph{看见}扰动的规划预见到这股推力、主动留裕度，真实峰值稳在墙内\cite{williams2018robust}。道理普适：\emph{rollout 必须反映真实系统会发生什么，否则算出的最优解只是对一个不存在的世界最优}。但真实扰动往往未知、时变，不能当已知常量塞进模型——Tube-MPPI/RMPPI 给的是处理\emph{未知}扰动的\emph{架构性}方案。

\paragraph{历史：从鲁棒 MPC 的 tube 到采样框架。}
鲁棒 MPC 有一经典思路 \textbf{tube（管道）}：把真实轨迹分解为“名义（无扰动）轨迹 $+$ 被反馈约束在管道内的误差”，只要误差始终留在管道里，名义安全即真实安全。Williams 等 $2018$ 年把它搬进 MPPI，提出 \textbf{Tube-MPPI}\cite{williams2018robust}；Gandhi 等 $2021$ 年指出其两个结构性缺陷并修复，提出 \textbf{RMPPI}，给出 MPPI 家族第一个性能保证\cite{gandhi2021robust}。值得注意：传统 tube 依赖\emph{线性误差动力学 $+$ 凸优化}（用不变集严格刻画管道），而 MPPI 的用武之地恰是\emph{非光滑代价、黑箱/非线性动力学}，那套“线性 $+$ 凸”的严格机器用不上。所以移植不是简单拼接——Tube-MPPI 保留“名义 $+$ 管道反馈”的\emph{架构思想}、放弃传统 tube 的硬集合保证，这也埋下它两个缺陷的根，以及 RMPPI 改用“自由能增长上界”这种更适配采样框架的保证形式的伏笔。

\begin{definition}[Tube-MPPI 双层架构]\label{def:mv-tube}
Tube-MPPI 把控制拆为两层：
\begin{itemize}
  \item \textbf{上层（名义控制器，nominal MPPI）}：在\emph{无扰动}的标称动力学上跑标准 MPPI，规划名义轨迹 $\bar{x}(\cdot)$ 与名义控制 $u_{\text{nom}}(\cdot)$；它须是采样式的，以处理稀疏/非光滑代价；用\emph{长}时域、\emph{低}频重规划。
  \item \textbf{下层（辅助控制器，ancillary，常用 iLQG）}：解一个\emph{跟踪}问题，把真实状态 $x$ 拉回名义轨迹 $\bar{x}$，使真实状态留在以 $\bar{x}$ 为中心的管道内；用\emph{短}时域、\emph{高}频运行，给出时变线性反馈 $u_t=\bar u_t+K_t(x_t-\bar x_t)$。
\end{itemize}
\end{definition}

\noindent 分工的道理：\emph{采样擅长全局、非光滑的规划，但不擅长高频纠偏；线性反馈擅长高频纠偏，但处理不了非光滑代价}。上下层的时域/频率\emph{不对称}（上层长而低频、下层短而高频）恰把“贵”（长时域全局规划）和“快”（短时域线性反馈）错开，合起来既远视又抗扰。下层增益由 iLQG（iterative Linear-Quadratic-Gaussian，iLQR 的高斯噪声版本）给出：在名义轨迹附近把动力学线性化（$A_t=\partial f/\partial x$，$B_t=\partial f/\partial u$）、代价二次化，得局部 LQ 问题，经\emph{反向 Riccati 递推}逐时刻算出一组时变反馈增益 $K_t$\cite{tassa2012ilqg}。

\begin{insight}
这里出现了贯穿现代机器人控制的范式——\textbf{“慢规划 $+$ 快线性反馈”的分层}。Riccati 增益把“如何在轨迹附近最优纠偏”压缩成一组可低频算好、在线高频查用的矩阵 $K_t$：纠偏只需一次矩阵乘 $K_t(x_t-\bar x_t)$，不必重解优化。Tube-MPPI/RMPPI 的双层正是它在采样式 MPC 里的体现——上层 MPPI 负责非凸、远视的规划，下层 Riccati 反馈负责局部、高频的纠偏，各取所长。
\end{insight}

\paragraph{planner--tracker 冲突：Tube-MPPI 的两个根因。}
Tube-MPPI 能抗扰，但它的两层\emph{各想各的}，由此带来两个结构性缺陷\cite{gandhi2021robust}：
\begin{enumerate}
  \item \textbf{名义状态与真实状态会发散。} 上层规划名义轨迹时，$\bar{x}$ 的演化\emph{独立于}真实状态 $x$ 选取；扰动持续把 $x$ 推离 $\bar{x}$，上层并不知情，两者越拉越远——状态发散是 Tube-MPPI 的主要失效模式。
  \item \textbf{采样器无视下层反馈。} 上层在围绕 $u_{\text{nom}}$ 撒噪声、评估 rollout 时，\emph{完全不知道}下层还有一个反馈在纠偏；于是采样出的轨迹反映的是“没有反馈时”的行为，给出对最优控制\emph{有偏}的估计。
\end{enumerate}
两缺陷共根：\emph{上层的提议分布建立在一个不完整的系统模型上}——既不看真实状态，也不看反馈的存在。换言之，Tube-MPPI 只是把反馈“挂”在 MPPI 外面，没让 MPPI 在采样时把反馈算进去。

\paragraph{RMPPI：把反馈嵌进重要性采样。}
RMPPI 的修法正是堵上这两个根，做两个关键改动\cite{gandhi2021robust}：

\emph{改动一——增广状态，让采样从真实状态出发。} 用一个\emph{增广状态空间}同时表示名义与真实动力学（把 $\bar x$ 与 $x$ 并成一个扩展向量，维度约翻倍），并配一套\emph{安全逻辑}：当名义与真实偏离过大时，把名义状态拉回真实状态附近，从源头避免发散。只有把两者放进\emph{同一个}状态空间，规划过程才能“同时看到”名义与真实、据二者之差决策——把“计划”与“实况”画在同一张图上，背离一目了然。

\emph{改动二——让 rollout 含反馈，权重却用名义代价。} 采样时每条 rollout 用的控制不再是“名义 $+$ 噪声”，而是
\begin{equation}
  u^\star_k(t)=u_{\text{nom}}(t)+K_{\text{fb}}\big(x_k(t)-\bar{x}(t)\big)+\varepsilon_k(t),
  \label{eq:mv-rmppi-rollout}
\end{equation}
其中 $K_{\text{fb}}$ 是下层反馈增益、$x_k(t)$ 是第 $k$ 条 rollout 的（含扰动的真实）状态、$\bar x(t)$ 是名义状态；rollout 用\emph{含扰动的真实动力学}前向，于是每条采样轨迹都“带着反馈”演化、反映真实闭环行为（堵缺陷二）。而权重仍用\emph{名义代价}$S_k$。

\begin{insight}
为何“rollout 用真实动力学、权重却用名义代价”看似拧巴却恰当？\emph{真实动力学回答“这条控制实际会把系统带到哪”}（要真实，堵缺陷二）；\emph{名义代价回答“这个去向对任务好不好”}（按名义任务评判）。扰动是要靠反馈\emph{吸收}的干扰项、而非优化目标本身；若权重也用含扰动的真实代价，优化目标就被随机扰动污染、每次“好坏”评判都带噪声，反而不稳。这正是 RMPPI 能在保持任务导向的同时获得鲁棒性的关键设计。回到 \cref{ch:mppi} 那条主线：MPPI 的一切改进都在让提议分布更接近真实最优分布，而在有扰动、有反馈的系统里，“真实”指的就是\emph{闭环}最优分布——RMPPI 让提议分布本身\emph{在闭环意义下}有效。
\end{insight}

\begin{theorem}[RMPPI 的自由能增长上界]\label{thm:mv-freenergy}
利用增广状态、安全逻辑与“考虑跟踪能力”的重要性采样，可推出系统\emph{自由能增长}的一个上界\cite{gandhi2021robust}。回顾 \cref{ch:mppi}：自由能刻画“当前提议分布离最优有多远”。一个对自由能\emph{增长}的上界，本质是在说“扰动作用下，系统离最优的偏离不会无限放大、被控制在可量化范围内”——正是鲁棒性的数学表达。该上界是三样东西的函数：\emph{任务约束的满足程度、底层跟踪控制器的性能、内部随机优化的采样误差}（直接对应 \cref{ch:mppi} 的 ESS）。三项任一变好，鲁棒性都可量化地提升。
\end{theorem}

\noindent 为何界的对象是自由能的\emph{增长}而非\emph{值}？在持续扰动的闭环系统里，逐周期谈“绝对的离最优距离”意义不大（每周期扰动在变、最优本身也在动）；真正决定系统会不会“越跑越偏、最终失控”的，是这个量\emph{每一步会不会被放大}——若每周期的增长都被一个界压住，误差就不会无限累积、系统稳定在最优附近的可控邻域内。这与控制理论用 Lyapunov 函数的\emph{差分}（而非函数值本身）证稳定性一脉相承。论文用两种跟踪器（iLQG 与基于收缩度量的控制器）验证了这个界，说明它\emph{不依赖}具体跟踪器；并在 GT AutoRally 越野车上演示 RMPPI 同时缓解“MPPI 缺乏鲁棒性”与“Tube-MPPI 过度保守”两个问题\cite{gandhi2021robust}。

\begin{pitfall}[思维陷阱：把“自由能上界”当成硬约束安全保证]
把“自由能增长有上界”理解成“约束一定不被违反、安全已被证明”。\textbf{为何错}：自由能上界是关于“离最优的偏离被控制在范围内”的概率/能量意义的保证，\emph{不是}传统鲁棒 MPC 那种“有界扰动下硬约束必然满足”的严格保证；MPPI 的约束本质上仍是软的（\cref{ch:mppi}）。\textbf{正解}：理解为“显著更鲁棒、且鲁棒性可量化”，而非“绝对安全”；安全攸关时仍叠加硬安全层（如控制屏障函数 CBF 滤波，见章末\nameref{sec:mv-frontier}的 Shield-MPPI）。
\end{pitfall}

\paragraph{何时 Tube 够、何时需 RMPPI。}
粗略判据看\emph{扰动相对反馈纠偏能力有多强}：若扰动温和、下层反馈能轻松把真实状态摁在名义附近（偏离始终很小），Tube-MPPI 两个缺陷都不严重、且实现更简单；若扰动强到让真实状态频繁大幅偏离（高速越野、强打滑），名义/真实发散成主要失效模式、采样偏差也变大，RMPPI 的两个改动才真正物有所值。另一考量是\emph{是否需要可量化的鲁棒性保证}：要把鲁棒写进设计指标或认证，就需要 RMPPI 的自由能上界；只求工程能用，Tube 往往已足够。

\subsection{Smooth-MPPI：input-lifting 从源头消除抖动}
\label{sec:mv-smooth}

\paragraph{动机：抖动是高斯采样的固有产物，不是 bug。}
把 vanilla MPPI 的控制输出画出来会看到\emph{高频颤动}——相邻时刻的控制值上下乱跳。根源在于：vanilla MPPI 在\emph{控制空间}里、对\emph{每个时刻独立}地撒高斯噪声 $\varepsilon_k(t)\sim\mathcal{N}(0,\Sigma)$。“每个时刻独立”意味着 $\varepsilon_k(t)$ 与 $\varepsilon_k(t{+}1)$ 毫无相关，加权平均后这种逐点独立的白噪声残留为高频分量、无法消除。\cref{ch:mppi} 用 Savitzky--Golay（SGF）\emph{事后}磨平，但那有两个本质问题：其一，它在最优解算完之后再磨，\emph{破坏了信息论最优性}（磨平后的控制不再是加权平均给出的最优控制）；其二，事后滤波是\emph{因果}滤波、引入\emph{相位滞后}，在快速机动时损害性能。更警惕的是：抖动在环境快速变化时会加剧，极端情况甚至让 MPPI 发散\cite{kim2022smooth}。

\paragraph{核心思想：在控制导数空间采样。}
Kim 等 $2022$ 年提出 \textbf{Smooth-MPPI}（论文题 \emph{Smooth MPPI Control without Smoothing}，“without smoothing”强调不靠事后平滑）\cite{kim2022smooth}。核心叫 \textbf{input-lifting}：既然在控制空间撒白噪声必带高频，那就改在\emph{控制导数空间} $\delta u=\mathrm{d}u/\mathrm{d}t$ 里撒噪声，再积分回控制：
\begin{equation}
  \delta u_k(t)\sim\mathcal{N}(0,\Sigma_\delta),\qquad u_k(t)=u_k(t{-}1)+\delta u_k(t)\cdot\mathrm{d}t.
  \label{eq:mv-inputlift}
\end{equation}
为何这就平滑了？因为控制 $u_k(t)$ 现在是导数噪声的\emph{累积和}（积分），而对白噪声做积分相当于一个\emph{低通滤波}——高频被积分天然压制，相邻时刻的控制因共享前面累积的历史而彼此相关、变化平缓。

\begin{theorem}[input-lifting 不破坏最优性]\label{thm:mv-lift}
input-lifting 的数学本质是\emph{扩展状态空间}\cite{kim2022smooth}。原系统 $x_{t+1}=f(x_t,u_t)$、控制 $u$；定义增广状态 $\tilde{x}=(x,u)$，把控制换成导数 $\delta u$：
\begin{equation}
  \tilde{x}_{t+1}=\begin{pmatrix}x_{t+1}\\ u_{t+1}\end{pmatrix}
  =\begin{pmatrix}f(x_t,\,u_t)\\ u_t+\delta u_t\cdot\mathrm{d}t\end{pmatrix},\qquad \text{新控制}=\delta u_t.
  \label{eq:mv-lift-aug}
\end{equation}
原控制 $u$ 现在是增广状态的一个分量、按一个积分器演化，真正被采样优化的“控制”是导数 $\delta u$。于是：\emph{平滑是结构内生的}（$u$ 由 $\delta u$ 积分而来，平滑写在增广动力学里，不是事后加的）；\emph{最优性不受影响}（在增广系统 $(\tilde{x},\delta u)$ 上，MPPI 是标准跑法，\cref{ch:mppi} 的自由能--KL 推导\emph{原样成立}）。
\end{theorem}

\noindent 这就解释了 SMPPI 把“控制变化率惩罚 $\ell_{\text{smooth}}=\alpha\|\mathrm{d}u/\mathrm{d}t\|^2$ 写进\emph{状态}代价”——$u$ 现在是状态，惩罚它的变化（即 $\delta u$ 的大小）当然是状态代价的事，而非控制代价（\cref{ch:mppi} 强调控制代价由重要性采样\emph{内在}决定、不可随意添加）。input-lifting 本质上是给系统\emph{串联了一个积分器}（串联积分器是控制理论改变相对阶、平滑输入、消除稳态误差的经典手段，PID 的 I 项即此），可\emph{选择提升阶数}：提升到二阶（采样 jerk、积分两次）会让控制的一阶导也连续、更平滑，但响应更迟钝——“提升到几阶”本身是按任务权衡的设计选择，而非越高越好。

\begin{insight}
SMPPI 把“抖动”拆成两个正交的轴。考虑控制序列这张二维表：横轴是一条 rollout 内部的时刻 $t$，纵轴是控制周期的迭代 $i$。\textbf{i-轴抖动}是相邻周期算出的序列之间的跳变——vanilla MPPI 本来就能压（\cref{ch:mppi} 的 warm-start 让相邻周期解平滑衔接）；\textbf{t-轴抖动}是\emph{同一条 rollout 内部}相邻时刻控制值之间的跳变——这正是 vanilla MPPI \emph{压不住}的那种（控制空间逐点独立白噪声）。Smooth-MPPI 真正治的是 \textbf{t-轴抖动}：input-lifting 把控制变成导数的累积积分，让同一条 rollout 内相邻时刻天然相关。这解释了为什么单纯的 warm-start（只管 i-轴）治不了抖动——抖动主要长在 rollout 内部的 t-轴上。
\end{insight}

\paragraph{什么时候不需要 Smooth-MPPI。}
三种情形可省：执行器本身就是低通的（带大惯量的机械臂、或控制先过硬件滤波）——高频指令到不了执行器，抖动无害；任务本就需要 bang-bang 式剧烈切换（某些时间最优控制的最优解就在边界间快速跳变，强行平滑反偏离最优）；纯仿真/离线规划不上真机（抖动只在真实执行器上造成磨损）。SMPPI 还提醒：input-lifting 要发挥敏捷优势，依赖\emph{足够准的系统辨识}——模型差时，平滑的好处会被模型误差吃掉。

\subsection{Log-MPPI：重尾采样保住杂乱环境的可行性下限}
\label{sec:mv-log}

\paragraph{动机：杂乱环境里所有 rollout 都撞了。}
前两小节修“提议中心在扰动下失效”（\cref{sec:mv-robust}）与“提议时间结构带来抖动”（\cref{sec:mv-smooth}）。这一小节面对更基础的困境：在\emph{高度杂乱}的环境（密集障碍、窄通道）里，从当前控制序列周围撒高斯噪声，绝大多数 rollout 一头撞进障碍、拿到极高代价。当\emph{几乎所有}$S_k$ 都因碰撞而巨大时，权重 $w_k\propto e^{-S_k/\lambda}$ 全部趋零、彼此难分高下，加权平均退化成对一堆“全是坏的”轨迹求平均，MPPI 输出近乎随机。这不是样本\emph{效率}问题（\cref{ch:mppi} 那条主线讲“有用样本占比”，前提是至少有一些可行样本被稀释），而是更尖锐的\emph{可行性}问题：你手里可能\emph{一条可行 rollout 都没有}。

\begin{note}
\textbf{MPPI 健康度的三个层次.}\;层次一：有大量可行样本（理想，vanilla 就好用）；层次二：可行样本稀少、被淹没（样本效率问题，warm-start/调 $\lambda$ 能救）；层次三：\emph{几乎没有}可行样本（可行性问题）。重尾采样针对层次三——先靠重尾的偶发大动作“捞”出哪怕几条可行的，把 MPPI 从层次三拉回层次二，让它重新有米下锅。所以杂乱环境卡死时应先查“有没有可行样本”（可行性），而非急着调 $\lambda$ 提“效率”——后者在层次三根本不适用。
\end{note}

\paragraph{核心思想：Normal-Log-Normal 重尾混合分布。}
Mohamed 等 $2022$ 年提出 \textbf{Log-MPPI}\cite{mohamed2022logmppi}：既然高斯钟形尾巴太薄、产生不了大动作，就给采样分布接一条重尾。具体用 \textbf{NLN（Normal-Log-Normal）混合分布}：
\begin{equation}
  \varepsilon \sim \pi\cdot\mathcal{N}(0,\sigma^2)\;+\;(1-\pi)\cdot\mathrm{LogN}(\mu_{\ln},\sigma_{\ln}^2),
  \label{eq:mv-nln}
\end{equation}
以概率 $\pi$ 从普通高斯采样（负责均值附近的精细探索 exploitation），以概率 $1-\pi$ 从对数正态采样（对称化后提供\emph{重尾}、偶发大动作 exploration）。这里“对称化”是必须做对的细节：对数正态天生\emph{单侧}（只取正值），而控制扰动必须\emph{双向}（既能加大也能减小油门、既能左转也能右转）；常见做法是给对数正态采样值随机乘 $\pm1$，让重尾对称分布在零两侧。漏掉对称化会导致隐蔽 bug——控制只在单方向探索得开、反方向被卡住，且因不报错极难调试。

\begin{insight}
为何不直接用\emph{均匀分布}（让大动作和小动作等概率）？因为均匀分布把概率均摊到整个范围，\emph{精细搜索的样本被稀释}——为偶尔的大动作牺牲了大量本该用于精细调优的样本。NLN 的高明正在\emph{两者兼顾}：主体仍是高斯（保留精细探索），只在尾部叠加重尾（提供偶发大动作）。这是“大部分精细 $+$ 少量大胆”的最优配比。所以重尾不是“让采样更狂野”，而是“在保留精细探索的前提下，额外开一个大胆探索的低概率通道”。注意它和宽高斯的形状差异：重尾是“尖峰 $+$ 厚尾”（多数样本集中、少数极远），宽高斯是“整体摊平”（处处稀疏）——对应不同的探索行为。
\end{insight}

\noindent 选\emph{对数正态}而非 Student-$t$/Cauchy 的考量：它与高斯有天然数学联系（取对数即高斯）、便于组合成解析可处理的混合；重尾程度可由 $\sigma_{\ln}$ 平滑调节（不像 Cauchy 那样尾巴重到连均值/方差都不存在）。关键不在于非得是对数正态，而在于\emph{用一个尾巴比高斯厚、但又不至于失控的分布}替换薄尾高斯。后续工作 \textbf{GP-MPPI}（IROS 2023）用稀疏高斯过程在线推荐子目标，给重尾探索装上“方向盘”——让大动作不只盲目地大、而是朝有希望的方向大\cite{mohamed2023gpmppi}。

\begin{pitfall}[思维陷阱：可行性问题总能靠“采样更猛”解决]
遇 MPPI 在杂乱环境卡死，一律靠加重尾、加大探索来救。\textbf{为何错}：极端杂乱（如必须走一条特定长绕行路径）时，纯靠\emph{无方向}的随机重尾命中可行路径的概率仍极低；一味加猛探索还让控制更抖、更不安全。\textbf{正解}：重尾解决“敢不敢迈大步”，方向问题要靠\emph{引导}（GP-MPPI 子目标、或上层规划器给路径）。把 Log-MPPI 当成“提高可行样本\emph{下限}”的手段，而非杂乱环境的万能解。
\end{pitfall}

\subsection{SVG-MPPI：Stein 变分引导的寻峰}
\label{sec:mv-svg}

\paragraph{动机：两条都安全的路，平均一下就撞墙。}
前三小节改的是提议分布的位置、时间结构、尾部，但都默认提议\emph{单峰}。可现实里最优动作分布常\emph{多峰}：绕过一个障碍，向左绕和向右绕可能同样好，对应两条分开的最优路径、是最优分布的两个峰（mode）。vanilla MPPI 只能用一个高斯近似这个双峰分布，而单峰高斯要同时覆盖两个分得很开的峰，其均值只能落在两峰\emph{中间}——那恰是障碍所在。

\begin{insight}
背后有个简单几何直觉：\textbf{逐元素平均是一种凸组合，而可行区域（绕开障碍的路径集合）不是凸集}。两条绕障路径分居障碍两侧，它们的“中点”自然落在障碍上——就像圆环上两个对称点的连线中点落在圆心。MPPI 的加权平均本质就是在求一个（加权）凸组合，当最优解集非凸（多个分离的可行区）时，凸组合可能落到禁区。这不是 MPPI 实现的瑕疵，而是“用一个点（加权均值）代表一个非凸集合”的固有矛盾。这是 vanilla MPPI 的第四个核心缺陷：\textbf{单模态}。
\end{insight}

\paragraph{为何调参治不了它。}
把温度 $\lambda$ 调小、权重更尖，在两 mode 代价\emph{有明显高低}时管用（权重集中到更优的峰）；但两 mode \emph{几乎等优}时（左右绕一样好），权重难分伯仲，$\lambda$ 再小也压不出唯一赢家，平均仍发生，且 $\lambda$ 过小会让 ESS 崩、退化成贪婪。更根本的是：单个高斯提议在\emph{结构上}就是单峰的，无论怎么调 $\lambda$/$\Sigma$ 都无法表达“要么往左、要么往右”这种\emph{离散选择}。这是分布的\emph{表达能力}问题，不是参数问题——必须改提议分布的\emph{形式}。

\paragraph{核心思想：用 Stein 变分侦察、锁定一个目标 mode。}
Honda 等 $2024$ 年提出 \textbf{SVG-MPPI}（Stein Variational Guided MPPI），专门处理快速机动车辆那种\emph{快速漂移}的多模态最优分布\cite{honda2024svgmppi}。关键思路有点反直觉：不是“同时覆盖所有 mode”，而是\emph{找出一个目标 mode 并引导解收敛到它}——论文称之为 \textbf{mode-seeking（寻峰）}。为何不覆盖所有峰？因为覆盖多峰会牺牲 MPPI 最宝贵的性质——\emph{闭式、不需迭代的快速收敛}。SVG-MPPI 用改进的 Stein 变分梯度下降（SVGD）\emph{粗略估计}出一个目标 mode，再嵌入 MPPI 求一个\emph{只覆盖该 mode}的闭式解，既避开“取两峰中间”的灾难，又保住快速收敛。

\begin{definition}[Stein 变分梯度下降]\label{def:mv-svgd}
SVGD 是一种确定性变分推断：维护一组粒子 $\{\theta_1,\dots,\theta_M\}$（每个 $\theta_i$ 是一条候选控制序列），迭代更新让其整体逼近目标分布 $p(\theta)\propto e^{-J(\theta)/\lambda}$：
\begin{equation}
  \theta_i \leftarrow \theta_i+\epsilon\cdot\frac{1}{M}\sum_{j=1}^{M}\Big[\,k(\theta_j,\theta_i)\,\nabla_{\theta_j}\log p(\theta_j)\;+\;\nabla_{\theta_j}k(\theta_j,\theta_i)\,\Big],
  \label{eq:mv-svgd}
\end{equation}
$k$ 为 RBF 核。\textbf{第一项（吸引/exploitation）}$k(\theta_j,\theta_i)\nabla\log p(\theta_j)$ 把粒子推向高概率（低代价）区，$\nabla\log p$ 指向代价下降最快方向；\textbf{第二项（排斥/exploration）}$\nabla_{\theta_j}k(\theta_j,\theta_i)$ 把彼此靠太近的粒子推开、防止坍缩到同一点。正是排斥项让粒子\emph{铺开}、感知到不同 mode 的存在——这是 SVGD 相对单个高斯的根本优势：用\emph{一群}粒子表示分布，自然能表达多峰。
\end{definition}

\begin{algo}[SVG-MPPI 四步流程]\label{alg:mv-svg}
用一个\emph{改进的}（轻量化、少粒子、少步数）SVGD 侦察一个好的目标 mode，再交给 MPPI 收敛\cite{honda2024svgmppi}：
\begin{enumerate}
  \item \textbf{输运 guide particles}：用改进 SVGD 最小化\emph{反向 KL}（RKL）散度，把一小组 guide particles 输运到分布高概率区；其\emph{输运轨迹}用来粗略估计一个收敛的目标 mode。
  \item \textbf{取峰为名义序列}：把粒子里\emph{序列代价最低}的那个直接当目标 mode 的峰、即 MPPI 的名义控制序列（\emph{不取均值}——取均值会重蹈“多模平均撞墙”）。
  \item \textbf{拟合协方差}：用高斯拟合结合粒子输运轨迹，粗略估计该 mode 的协方差，得自适应 $\Sigma$。
  \item \textbf{mode-seeking MPPI}：以上面估出的均值与协方差跑一次标准 MPPI 采样——均值/协方差已锁定单个 mode，这次采样只覆盖该 mode，加权平均不再跨峰。
\end{enumerate}
\end{algo}

\noindent 第 1 步用\emph{反向 KL} 是 mode-seeking 的数学根源。逼近一个分布 $p$ 有两种 KL 方向：\emph{前向 KL} $\mathrm{KL}(p\,\|\,q)$ 是 mode-covering（覆盖型）——惩罚“$p$ 有概率而 $q$ 没有”，逼 $q$ 盖住所有峰，单高斯逼近双峰时 $q$ 摊在两峰之间（正是那个撞墙的平均，对应矩匹配）；\emph{反向 KL} $\mathrm{KL}(q\,\|\,p)$ 是 mode-seeking（寻峰型）——惩罚“$q$ 有概率而 $p$ 没有”（零强制 zero-forcing），逼 $q$ 缩进\emph{某一个}峰、不敢跨到峰间低概率区。vanilla MPPI 的加权平均本质是矩匹配（用样本估计目标分布的均值），所以继承前向 KL 的“摊到中间”毛病；SVG-MPPI 改用 RKL，正为换上“只占一峰”的行为。

\begin{insight}
SVG-MPPI 的精髓是“用 SVGD 侦察、用 MPPI 收敛”的分工——SVGD 擅长感知多峰但迭代慢，MPPI 闭式快但只有单峰，让 SVGD 只干“挑哪个峰”这件轻量活、把“在选定峰内快速收敛”交给 MPPI。这与另一条更“正统”的多模路线 \textbf{SV-MPC}\cite{lambert2020svmpc}（用 SVGD 直接输运一大群样本逼近\emph{整个}多峰后验）形成对照：SV-MPC 覆盖多峰但\emph{求不出闭式解、必须迭代}，迭代会劣化收敛、限制并行度与每次可用样本数，丢掉 MPPI 闭式快速与 GPU 并行的优势。一句话：\emph{SV-MPC 用 SVGD 扛起整个推断（慢、样本少、覆盖多峰）；SVG-MPPI 用 SVGD 只做轻量侦察、重活交给闭式 MPPI（快、样本多、只锁一个 mode）}。mode-seeking 是一种\emph{有意的偏置}——用“放弃覆盖其他 mode”换“闭式快速收敛”；任务只需一条好路径时这交易划算，需同时保留多个候选（如随时在左/右绕间切换以应对动态障碍）时则需真正多模的方法。
\end{insight}

\subsection{U-MPPI 与 CoVO-MPC：不确定性传播与最优协方差}
\label{sec:mv-umppi}

\paragraph{动机：MPPI 看得见“撞没撞”，看不见“多大概率撞”。}
前四小节修的都是\emph{提议分布}本身（位置、时间结构、尾部、模态）。这一小节换角度：\emph{rollout 怎么算}。vanilla MPPI 的每条 rollout 是一条\emph{确定}轨迹——给定一个噪声样本，前向积分出一条状态序列，看它撞没撞。但真实系统有不确定性（状态估计噪声、模型误差、扰动），系统从某状态出发下一步其实是一个\emph{分布}而非一个点。vanilla MPPI 不显式传播这个分布，于是只能回答“这条采样轨迹撞没撞”，回答不了“考虑不确定性，从这里走有多大概率撞”。一条擦着障碍边缘没撞的轨迹（代价低）可能其实很危险——只要状态稍偏就撞，碰撞\emph{概率}很高。

\begin{note}
\textbf{两种不确定性，别和 \cref{sec:mv-robust} 混.}\;其一是\emph{模型失配/扰动}——真实演化偏离标称模型（\cref{sec:mv-robust} 的 RMPPI 处理这个，靠反馈把状态拉回）；其二是\emph{状态估计噪声/过程噪声}——对当前状态本身就测不准、或动力学带随机性（本小节处理这个，靠传播量化）。一个用反馈\emph{对抗}不确定性，一个用传播\emph{量化}不确定性——视角互补，也能叠加。
\end{note}

\paragraph{核心思想：用少数 sigma 点高效传播均值与协方差。}
Mohamed 等 $2025$ 年提出 \textbf{U-MPPI}\cite{mohamed2025umppi}，在 MPPI 里嵌入 \textbf{Unscented Transform（UT）}：不用大量随机样本，而用 $2n_x+1$ 个\emph{确定性}的 sigma 点传播状态的均值和协方差（$n_x$ 为状态维数）。UT 源自 Julier 与 Uhlmann 为非线性滤波提出的工具（最知名应用是无迹卡尔曼滤波 UKF\cite{julier1997unscented}）；U-MPPI 把它从\emph{滤波/估计}迁移到\emph{控制/规划}的 rollout 里。UT 三步：\emph{生成 sigma 点}（围绕均值 $\mu_x$、按协方差 $P$ 摆 $2n+1$ 个确定性点，中心一个、沿协方差矩阵\emph{平方根}的每个方向正负各一）
\begin{equation}
  \chi_0=\mu_x,\quad \chi_i=\mu_x+\big(\sqrt{(n{+}\kappa)P}\big)_i,\quad \chi_{i+n}=\mu_x-\big(\sqrt{(n{+}\kappa)P}\big)_i,\quad i=1\dots n;
  \label{eq:mv-sigma}
\end{equation}
\emph{传播}（每个 sigma 点过非线性动力学一步 $\chi_i'=F(\chi_i,u+\varepsilon)$）；\emph{加权重组}
\begin{equation}
  \mu_x'=\sum_i w_i^m\,\chi_i',\qquad P'=\sum_i w_i^c\,(\chi_i'-\mu_x')(\chi_i'-\mu_x')^\top.
  \label{eq:mv-ut-recombine}
\end{equation}
对高斯（及近高斯）分布，UT 能以 $2n+1$ 个点精确捕获到均值与协方差（直至三阶矩），比蒙特卡洛高效几个数量级，且\emph{不需算雅可比}（对黑箱/不可微动力学也能用，与 MPPI 的黑箱友好一脉相承）——这是“efficient”的来源。

\begin{insight}
有了高效传播的均值和协方差，U-MPPI 能做 vanilla 做不到的两件事。其一，\emph{风险敏感的代价}：评估每条 rollout 时把“这条路有多不确定”算进去（论文报告它在室内杂乱环境的碰撞率显著低于 vanilla 与 Log-MPPI，且\emph{增强探索、降低陷局部极小的风险}——风险信息本身成了引导探索的线索）。其二，\emph{机会约束（chance constraint）}：其扩展 C²U-MPPI 用传播出的协方差\emph{估算碰撞概率}，当某轨迹 $P(\text{collision})>\delta$（用户设定的风险阈值）时加重罚——把约束从 vanilla 那种“撞了才罚”的硬碰硬，变成“快要太危险就提前规避”的概率约束。直觉上 UT 给出每时刻的均值 $\mu_x'$ 与协方差 $P'$（“大概在 $\mu_x'$、有 $P'$ 这么大的不确定椭球”），把椭球与障碍几何叠加即得碰撞概率——同样一条“擦边”轨迹，协方差大则概率高、被重罚，协方差小则概率低、罚得轻。这正是 vanilla 缺的一环：它只有均值（“大概在哪”），没有协方差（“可能偏多远”），分不清“稳稳擦过”与“侥幸擦过”。
\end{insight}

\paragraph{CoVO-MPC：把协方差从手调变成最优设计。}
六大变体里前五个改的是提议的某个侧面或 rollout 的刻画，但采样协方差 $\Sigma$ 大多仍\emph{手工设定}。\textbf{CoVO-MPC}（Yi 等 $2024$，L4DC）给出第六个角度——\emph{最优协方差设计}，并为采样式 MPC 的\emph{收敛性}给出理论刻画：当代价二次（涵盖时变 LQR）时 MPPI \emph{至少线性收敛}，并据此推出最优采样协方差的闭式\cite{yi2024covo}。其结论是：最优 $\Sigma$ 与代价 Hessian $D$ 共享特征向量，且在“采样体积 $\det\Sigma$ 受限”的约束下取 $\Sigma\propto D^{-1/2}$——即\emph{在代价陡峭（$D$ 大）的方向上收窄、在平缓（$D$ 小）的方向上放宽}。直觉是：vanilla 的各向同性 $\Sigma$ 不知道代价地形的“形状”，在陡峭方向撒太宽、在平缓方向撒太窄都低效；CoVO 按局部代价地形（Hessian）把 $\Sigma$ 整形到最优方向，让样本沿“最该探索”的方向铺开，从而更快收敛（实测在敏捷四旋翼控制上比标准 MPPI 优 $43$--$54\%$）。它和 U-MPPI 共享“把一个靠经验拍脑袋的量（裕度 / 协方差）换成由可观测信息驱动的自适应量”的工程思路——这也呼应 \cref{sec:mv-robust} 把“鲁棒”拆成可调旋钮、\cref{ch:mppi} 把“采样健康”变成 ESS 的同一种思维。

\subsection{统一视角与选型}
\label{sec:mv-select}

\paragraph{五侧面表：六个变体认领高斯的五个侧面。}
把全节反复出现的“侧面”观点正式收口。vanilla MPPI 用一个\emph{开环、各向同性、单峰、薄尾、只传均值}的高斯近似 $\mathbb{P}^\star$——这五个限定词正是它的五个短板，也正是六个变体各自要松开的那一个。

\begin{table}[H]\centering\small
\caption{提议分布的五个侧面与对应变体}
\label{tab:mv-sides}
\begin{tabular}{p{0.13\textwidth} p{0.26\textwidth} p{0.18\textwidth} p{0.30\textwidth}}
\toprule
侧面 & vanilla 的假设 vs 真实 $\mathbb{P}^\star$ & 变体 & 怎么松 \\
\midrule
位置（中心） & 围绕旧解、开环 vs 围绕真实状态的闭环解 & RMPPI（\cref{sec:mv-robust}） & 反馈把中心拉回真实状态、采样知道有反馈 \\
时间结构 & 各时刻独立白噪声 vs 时间相关、平滑 & Smooth-MPPI（\cref{sec:mv-smooth}） & 导数空间采样再积分（积分$=$低通） \\
尾部 & 薄尾 vs 重尾（偶有大动作） & Log-MPPI（\cref{sec:mv-log}） & NLN 重尾混合分布 \\
模态 & 单峰 vs 多峰（多条等优路径） & SVG-MPPI（\cref{sec:mv-svg}） & SVGD 侦察、锁定一个目标 mode \\
不确定性刻画 & 只传均值 vs 带协方差的轨迹分布 & U-MPPI（\cref{sec:mv-umppi}） & UT 传播均值 $+$ 协方差 \\
协方差形状 & 各向同性、手调 vs 贴合地形的最优形状 & CoVO-MPC（\cref{sec:mv-umppi}） & 按代价地形整形 $\Sigma$，带收敛理论 \\
\bottomrule
\end{tabular}
\end{table}

\begin{algo}[选型决策树]\label{alg:mv-decision}
按“哪个缺陷最致命”自上而下、命中第一个匹配分支即可（不是把所有分支都堆上）：
\end{algo}

\begin{center}
\small\textbf{选型决策树（按最致命的痛点）}\\[3pt]
\begin{forest} dtree
[{你的任务最大的\\痛点是什么?}
  [{扰动/模型失配导致\\约束被违反、行为发散?}
    [{→ RMPPI(要理论保证)\\\quad / Tube-MPPI(实现简单)}]
  ]
  [{控制抖动不可接受\\(执行器娇贵、需平滑)?}
    [{→ Smooth-MPPI}]
  ]
  [{环境高度杂乱、\\rollout 几乎全碰撞?}
    [{→ Log-MPPI\\\quad / GP-MPPI(带子目标引导)}]
  ]
  [{存在多条等优路径、\\怕取中间撞墙?}
    [{→ SVG-MPPI(寻峰)}]
  ]
  [{需要风险/不确定性感知、\\想控制碰撞概率?}
    [{→ U-MPPI / C2U-MPPI}]
  ]
  [{需要收敛保证、\\想最优地设计协方差?}
    [{→ CoVO-MPC}]
  ]
  [{以上都不突出、标称模型够准、\\环境不杂乱?}
    [{→ vanilla MPPI 就够了}]
  ]
]
\end{forest}
\end{center}

\begin{insight}
两条选型心法。其一，\textbf{别按论文新旧或“听起来先不先进”来选，而按缺陷选}：U-MPPI（$2025$）比 Tube-MPPI（$2018$）新，不代表“更好”——它们解的是完全不同的缺陷（不确定性 vs 模型失配），新的那个对你的任务可能毫不相关。其二，\textbf{变体大多可正交叠加}（改不同侧面）：Smooth（时间结构）$+$ Log（尾部）可直接相乘（在导数空间用 NLN 采样）；但也有内在张力需调平衡——Log 的重尾甩出的偶发大动作\emph{本身是高频跳变}，与 Smooth 的平滑目标对着干，叠加时重尾比例 $\pi$ 不能太低、导数空间协方差也要相应收一点。实操心法：\emph{先按需求优先级单独加变体、each 验证有效，再两两组合、观察指标是否如预期叠加}，不要一上来把所有开关全打开。
\end{insight}

\begin{practice}
\begin{enumerate}
  \item 写出 RMPPI 相对 Tube-MPPI 的两个改动，并说明各堵哪个缺陷。再论证：若把 \cref{eq:mv-rmppi-rollout} 的反馈项漏掉、只在执行端套反馈，为什么会“退化成换皮 Tube-MPPI”、自由能上界的前提为何失效。
  \item 用增广状态 $\tilde x=(x,u)$ 推导 input-lifting 后的动力学 \cref{eq:mv-lift-aug}，并解释为什么“控制变化率惩罚”应写进\emph{状态}代价而非控制代价（提示：回顾 \cref{ch:mppi} 控制代价由重要性采样内在决定）。
  \item 构造一个“向左绕、向右绕”等优的双峰避障场景，验证两条路各自安全、但逐元素平均后撞障碍；再实现一个最朴素的“寻峰”（按终点位置把样本聚成两簇、各自加权、取代价更低的簇输出），验证它能避障。把它和 \cref{alg:mv-svg} 的第 2 步对应起来。
\end{enumerate}
\end{practice}

% ════════════════════════════════════════════════════════════════
\section{扩散同构：MPPI 作为单步去噪，退火征服高维}
\label{sec:mv-diffusion}

\paragraph{动机：扩散模型凭什么生成高斯生不出的东西？}
\cref{sec:mv-six} 在“单步高斯框架”内把提议分布的每个侧面都推进了一步，但框架本身锁死了上限——单峰高斯逼近多峰目标会卡在两峰间的谷底（\cref{sec:mv-svg}）。本节换一个维度突破：扩散模型偏偏以“能生成极复杂、多峰分布”闻名，而它\emph{每一步用的恰恰也是高斯}（反向过程每步加高斯噪声、学高斯条件分布）。既然单步还是高斯，它凭什么生成高斯生不出的复杂分布？

\begin{insight}
答案是本节最重要的观念：\textbf{扩散模型的表达力不来自任何单步，而来自“许多步高斯去噪的复合”。} 一步高斯挪不远、雕不出复杂形状，但几十上百步首尾相接，就能把一团纯噪声逐步“雕刻”成任意复杂的多峰分布。这是“用深度换宽度”的智慧——神经网络用多层换表达力、扩散用多步换表达力，本质相通。它对采样式 MPC 的提示是：\emph{不必去设计一个更花哨的单步提议分布（那很难、不通用），只要学会把“许多步高斯更新”组织成一个去噪过程，就能获得远超单步高斯的表达力}。而你早就会“一步高斯更新”了——那就是 MPPI。
\end{insight}

\subsection{DDPM 与 MPPI 的数学桥梁}
\label{sec:mv-bridge}

\paragraph{历史：DDPM 与 score-based 两条线汇流。}
扩散模型是两条独立技术线汇流的产物。\textbf{DDPM}（去噪扩散概率模型，Ho 等 $2020$）\cite{ho2020ddpm}把生成组织成“前向逐步加噪、反向逐步去噪”，训一个网络预测每步该去掉的噪声（语言是“噪声预测”）；\textbf{score-based 生成}（Song--Ermon $2019$）\cite{song2019score}从另一角度切入——把生成理解为“沿数据分布的对数梯度（score）$\nabla_x\log p$ 往上爬”（语言是“score 估计”）。两条线很快被证明\emph{数学等价}（差一个已知系数，即下文的 Tweedie 关系），于是“去噪一步”和“朝 score 上升一步”是同一件事。控制借的主要是 score-based 这条线的语言（“朝低代价/高密度方向优化”与控制天然契合），而 DDPM 那条线贡献了“多级噪声调度”的框架（对应退火调度）。

\begin{theorem}[MPPI 一步 $=$ 单步去噪 $=$ 一步 score 上升]\label{thm:mv-isomorphism}
把 MPPI 更新 \cref{eq:mv-mppi-update} 与 DDPM 反向步并排，二者是同一个模板。DDPM 在去噪某一级，给定带噪估计 $x_k$，反向一步生成更干净的
\begin{equation}
  x_{k-1}=\underbrace{x_k+c_k\,\nabla_x\log p_k(x_k)}_{\text{朝 score 方向去噪}}+\underbrace{\sigma_k\,z}_{\text{注入少量噪声}},\qquad z\sim\mathcal{N}(0,I).
  \label{eq:mv-ddpm-reverse}
\end{equation}
MPPI 的“$\sum_k w_k\varepsilon_k$”和 DDPM 的“$c_k\nabla\log p_k$”占同一位置（都是“指向更优/更高概率方向的修正量”）。考虑 MPPI 隐含的目标分布 $p(U)\propto\exp(-J(U)/\lambda)$，其 score 为 $\nabla_U\log p(U)=-\tfrac{1}{\lambda}\nabla_U J(U)$；而 MPPI 的加权噪声 $\sum_k w_k\varepsilon_k$ 正是该 score（经高斯平滑后）的一个\emph{蒙特卡洛估计}——MPPI 没解析地算 $\nabla J$，而是用“撒扰动、按 $e^{-J/\lambda}$ 加权”这个\emph{无导数}办法估出了“朝 score 上升”的方向\cite{pan2024mbd}。
\end{theorem}

\begin{derivation}[采样方差就是平滑尺度]
\cref{thm:mv-isomorphism} 用凸的二次代价验证时余弦相似度极高，但本节通篇讲非凸——同构对非凸代价还成立吗？成立，但要把“成立的是什么”说精确：MPPI 的加权平均估计的\emph{不是}原代价 $J$ 的 score，而是“被高斯核平滑过的目标分布”的 score。在当前点 $U$ 附近撒方差 $\sigma^2$ 的高斯扰动、按 $e^{-J/\lambda}$ 加权平均，算的是目标分布 $p(U)\propto e^{-J/\lambda}$ 与一个方差 $\sigma^2$ 高斯核\emph{卷积}后那个分布 $p_\sigma$ 在当前点的 score：
\begin{equation}
  \sum_k w_k\varepsilon_k \;\approx\; \sigma^2\,\nabla_U\log p_\sigma(U),\qquad p_\sigma=p * \mathcal{N}(0,\sigma^2 I).
  \label{eq:mv-smoothed-score}
\end{equation}
对凸的二次代价，$p_\sigma$ 还是同型单峰，其 score 方向与 $-\nabla J$ 一致（故余弦$\approx1$）；对非凸代价，$p$ 多峰，卷积后的 $p_\sigma$ 是“被模糊过的多峰”（$\sigma$ 越大模糊越狠、局部峰被抹得越平）。这恰是退火能工作的根本——\textbf{采样方差 $\sigma$ 既是采样范围、又是平滑尺度，加权平均就是平滑后分布的 score 估计}。这也解释了为什么 \cref{thm:mv-isomorphism} 的 demo 里 MPPI 增量的\emph{大小}远小于 $-\nabla J$、只有\emph{方向}一致：增量是 $\sigma^2\nabla\log p_\sigma$，那个 $\sigma^2$ 与平滑都缩小幅度、但不改方向。
\end{derivation}

\begin{insight}
\cref{ch:mppi} 那个“MPPI $\leftrightarrow$ REINFORCE/回报加权回归”的同构，到这里有了第三张面孔：REINFORCE 的回报加权、MPPI 的代价加权、扩散的 score 去噪，本质是同一个“加权估计上升方向”的操作。$2025$ 年已有工作把 MPPI、策略梯度、扩散\emph{显式统一}在同一个 Gibbs 测度 $q(U)\propto\exp(\hat E(U)/\tau)$ 下\cite{li2025unified}（差别只在能量 $\hat E$ 是什么：负代价 / 期望回报 / 对数密度），共享同一更新模板“按指数化能量给样本加权、朝加权方向挪一步”。一个采样式控制器、一个策略梯度、一个生成模型的去噪步，在数学最内核处是同一个东西。
\end{insight}

\subsection{MBD：用模型直接算 score，免训练免数据}
\label{sec:mv-mbd}

\paragraph{动机：想要扩散的本事，又不想要它的数据和训练。}
DDPM 与 Diffusion Policy\cite{chi2023diffusion} 的 score 是\emph{从数据训练出来的网络}——这正是它们强大、也正是它们受限处：要大量高质量演示，且 score 网络绑定训练时的动力学，换个机器人就得重训。\textbf{MBD（Model-Based Diffusion，Pan 等 $2024$）}\cite{pan2024mbd} 提出一个反转，出发点是被忽视的事实：\emph{在轨迹优化里，动力学和代价通常是已知的}。既然目标分布由代价和动力学\emph{完全确定}，score 就是一个可以\emph{算}出来的量，何必从数据\emph{学}？

\begin{theorem}[MBD 把轨迹优化重写成反向扩散]\label{thm:mv-mbd}
轨迹优化要找让代价 $J(\tau)$ 最小的轨迹 $\tau$，改写为从目标分布采样：$p_0(\tau)\propto\exp(-J(\tau)/\lambda)$（\cref{ch:mppi} 的 $\mathbb{P}^\star$ 同一对象，$\lambda\to0$ 时质量全集中到最优轨迹）。仿 DDPM 定义前向加噪 $\tau_k=\sqrt{\bar\alpha_k}\,\tau_0+\sqrt{1-\bar\alpha_k}\,\varepsilon$（$\bar\alpha_k$ 随 $k$ 从 $1$ 递减到 $0$），反向去噪每一步估计后验均值 $\hat\tau_0(\tau_k)=\mathbb{E}[\tau_0\mid\tau_k]$ 并朝它挪一步。关键是这个后验均值\emph{由模型算出}：
\begin{equation}
  \hat\tau_0(\tau_k)=\mathbb{E}[\tau_0\mid\tau_k]=\frac{\sum_i \exp(-J(\tau_0^i)/\lambda)\,\tau_0^i}{\sum_i \exp(-J(\tau_0^i)/\lambda)},\qquad \tau_0^i\sim\mathcal{N}(\tau_k,\sigma_k^2 I).
  \label{eq:mv-mbd-posterior}
\end{equation}
在 $\tau_k$ 附近撒一批候选 $\tau_0^i$，用模型对每条算代价（一次前向 rollout），按 $e^{-J/\lambda}$ 加权平均——\textbf{这就是 MPPI 的加权平均}，只不过 MBD 把它解释成“对干净轨迹的后验均值估计 / 一步 score 上升”。权重里的代价来自\emph{已知}的动力学和代价模型，\emph{不来自任何训练好的网络、不需要任何数据}\cite{pan2024mbd}。
\end{theorem}

\begin{derivation}[Tweedie 恒等式：score 与后验均值是一回事]
“朝 score 去噪”与“朝后验均值挪”由 Tweedie 公式联系。设噪声级 $k$ 的带噪样本 $\tau_k$ 是干净样本加方差 $\sigma_k^2$ 高斯噪声，则
\begin{equation}
  \nabla_{\tau_k}\log p_k(\tau_k)=\frac{\hat\tau_0(\tau_k)-\tau_k}{\sigma_k^2},\qquad \hat\tau_0(\tau_k)=\mathbb{E}[\tau_0\mid\tau_k].
  \label{eq:mv-tweedie}
\end{equation}
左边是 score（方向），右边是“后验均值减当前位置”再除以噪声方差——也是方向。它说 score 恰指向后验均值 $\hat\tau_0$，故“朝 score 上升一步”与“朝 $\hat\tau_0$ 挪一步”是同一动作（差一步长 $\sigma_k^2$）。与 \cref{eq:mv-mbd-posterior} 一拼，整条逻辑链闭合：\emph{在 $\tau_k$ 附近撒候选、用模型算代价、加权平均得 $\hat\tau_0$（MPPI），而 $\hat\tau_0-\tau_k$ 正比于 score（Tweedie），所以挪向 $\hat\tau_0$ 就是朝 score 上升一步（去噪）}。这也解释了无导数优化为何能逼近梯度方向——后验均值的位移本身就编码了 score 信息。注意 Tweedie 要求“候选是高斯加噪”，而退火采样 MPC 正用 $\tau_0^i\sim\mathcal{N}(\tau_k,\sigma_k^2 I)$ 撒候选，构造上自洽——这也说明 \cref{ch:mppi} “单步用高斯”不只是历史选择，还与这套去噪解释的自洽性有关。
\end{derivation}

\begin{insight}
MBD 颠覆的不是扩散的机制，而是 score 的\emph{来源}——把“从数据学 score”换成“用模型算 score”。之所以可能，是因为在轨迹优化里目标分布 $p_0\propto e^{-J/\lambda}$ 被代价和动力学\emph{完全确定}，其 score 是确定、可估计的量，根本不是需要从数据猜的未知函数。“用扩散就要有数据、要训练”的成见由此被打破：\emph{模型就是最好的、无限量的、且天然跨场景的监督信号}。MBD 还能\emph{融合不完美数据}（把演示当目标分布的“观测”、叠到 score 估计里）——因为模型那一项保证可行性底线，数据只在底线之上做“风格微调”，所以数据可以不完美、不完整、甚至动力学不可行（如用下半身人类动作引导全身人形）。论文摘要明确报告 MBD 在一系列（富接触）任务上比 PPO 高 $59\%$，且只需几十秒的“扩散”时间\cite{pan2024mbd}（原文：\emph{“MBD outperforms PPO by 59\% in various tasks within tens of seconds of diffusing.”}）。
\end{insight}

\subsection{DIAL-MPC：多级退火征服全阶力矩控制}
\label{sec:mv-dial}

\paragraph{动机：高维力矩控制是单级 MPPI 的滑铁卢。}
\cref{ch:mppi} 警告：高维下随机样本几乎都落在远离均值的薄壳上，“撒满”空间所需样本数随维度指数爆炸。在几十维的\emph{全阶力矩}空间（直接在完整刚体动力学上优化、直接输出关节力矩、不依赖底层跟踪器）里，单级 MPPI 要找到一条好轨迹需要的样本数大到不切实际——这正是采样式 MPC 长期被认为“上不了全阶腿足控制”的原因。\textbf{DIAL-MPC}（Xue 等 $2025$，ICRA 最佳论文入围）的破解正是 \cref{sec:mv-mbd} 的扩散退火：\emph{把“一次在几十维里撒满”这个不可能任务，拆成“多级逐步聚焦”的可行任务}\cite{xue2025dial}。

\begin{algo}[DIAL-MPC 主干：多级退火 MPPI]\label{alg:mv-dial}
算法上就是“多级退火 MPPI”，针对实时控制做两点工程化：warm-start（用上周期解平移作初值，利用时间连续性）与 bi-level 退火。主干是一个多级退火循环，每一级是一次标准 MPPI 更新，级与级之间温度和采样方差\emph{逐级缩小}：
\end{algo}

\begin{codebox}[text]{DIAL-MPC 主干（一个控制周期）}
输入: 当前状态 x0, 上周期解 U_prev, 退火级数 L, 每级 {(lambda_l, Sigma_l)}
U <- shift(U_prev)                       # 1. warm-start: 上周期解平移一格做初值
for l = 1..L:                            # 2. 多级退火 (= 反向扩散的 L 步去噪)
    for k = 1..K (并行):                 #    2a. 在当前 U 附近采 K 条扰动
        eps_k ~ N(0, Sigma_l)
        J_k = rollout_full_order(x0, U + eps_k)   # 全阶动力学 rollout(GPU 并行)
    rho = min_k J_k                      #    2b. 标准 MPPI 加权更新 (= 一步去噪)
    w_k = exp(-(J_k - rho)/lambda_l) / sum_j exp(-(J_j - rho)/lambda_l)
    U <- U + sum_k w_k * eps_k
    # lambda_l, Sigma_l 已按调度逐级缩小(高温大方差 -> 低温小方差)
execute(U[0])                            # 3. 只执行第一格, 下周期重来(滚动时域)
\end{codebox}

\noindent 与 vanilla MPPI 的差异只有一处但关键：vanilla 在单个 $(\lambda,\Sigma)$ 下更新\emph{一次}（单步去噪），DIAL-MPC 用 $L$ 级、每级 $(\lambda_l,\Sigma_l)$ 逐级缩小（完整反向扩散）。\textbf{bi-level（双层）退火}是其精细之处：\emph{轨迹级退火}（trajectory-wise）在不同在线迭代上用不同采样方差、对整条轨迹逐步精炼；\emph{动作级退火}（action-wise）对规划时域里\emph{不同位置}的控制用不同退火——越靠近当前时刻越需精确（马上执行），越远的越可“粗”（后面还会重规划）。两层叠加，把有限采样预算花在刀刃上。

\begin{insight}
为什么退火能赢？它把“探索”与“利用”在\emph{时间上分开}：高温大方差阶段牺牲精度换全局视野（先看清深阱在哪），低温小方差阶段牺牲视野换精度（在深阱里收敛）。单级 MPPI 被迫用一个固定尺度同时干这两件相互矛盾的事，自然两头不讨好。这和 DDPM 反向过程“从高噪声级逐步去噪到低噪声级”是同一道理——\emph{高噪声级管全局、低噪声级管局部，逐级交棒}。用 \cref{eq:mv-smoothed-score} 的语言：高 $\sigma$ 级估计的是“重度模糊版”地形的 score（能看全局、抹平了小峰），低 $\sigma$ 级估计的是“轻度模糊版”的 score（能精修）——退火让平滑核从大到小，先抓大势、后抠细节，从而绕过困住单级方法的局部陷阱。这和图像处理的多分辨率金字塔、优化里的同伦延拓、人解决问题“先抓主要矛盾、再处理细节”是同一智慧的不同化身。
\end{insight}

\noindent DIAL-MPC 之所以能做全阶力矩控制，本质是两个“免”叠加：\emph{免求导}（采样式，绕开接触切换处的不可微动力学——脚落地/离地时接触力瞬间出现/消失，梯度类方法 iLQR/DDP 在此失效）与\emph{免陷局部}（多级退火做全局搜索，绕开高维强非凸的局部最优）。任一单独不够：只有采样、没有退火，在几十维上会陷局部；只有退火思想、没有采样而用梯度，接触一来就崩。论文报告：比标准 MPPI 跟踪误差低 $13.4\times$、接触奖励高 $107.7\%$；相比 GCRL 在\emph{行走}任务跟踪误差低 $3.9\times$；在高难\emph{爬越}（crate climbing）任务比 RL 策略高 $50\%$（且是唯一能稳定给出可行解的采样式 MPC，成功率比最优 MPPI 调度高 $3\times$）；并在真实 Unitree Go2（18 自由度、12 驱动）上以 $50$ Hz 完成带 $7$\,kg 负载的跳跃，全程\emph{免训练}\cite{xue2025dial}。（注：正式版报告的是“爬越任务比 RL 高 $50\%$”，并非早期流传的“跳跃高 $40\%$”。） 它的成立一半靠算法洞见（退火），另一半靠工程时机（GPU 大规模并行物理仿真，如 MuJoCo MJX/Brax）——这呼应 \cref{ch:mppi} “采样式方法的复兴被 GPU 并行能力催生”的故事，在全阶腿足这个更极端规模上再演一遍。

\begin{insight}
把全章主线收成一句：第 \cref{ch:mppi} 章留下高斯天花板（单步是高斯、单峰逼近多峰会卡谷底，且在单步框架内无法打破），本节的回答是\textbf{不在单步上较劲，改用多步的复合}——每步仍是朴素好采样的高斯（\cref{thm:mv-isomorphism}），但把许多步沿“由高到低的噪声/温度调度”复合起来（退火），就能自动生成任意复杂的多峰分布、爬上几十维全阶控制（\cref{alg:mv-dial}）；而多步去噪所需的“往哪走”（score），在动力学和代价已知时由模型直接算出、无需训练任何网络（\cref{thm:mv-mbd}）。所以“扩散给采样式控制带来了什么”的答案是：它没有发明一个更聪明的单步，而是证明了“把简单的单步复合起来”就能突破单步的天花板——\emph{威力从来不在某一步有多花哨，而在许多步如何被组织}。
\end{insight}

\begin{pitfall}[概念误区：“扩散启发的采样 MPC”就是要训练一个扩散网络]
一听“扩散 MPC”就默认要收集数据、训一个扩散模型生成动作（像 Diffusion Policy 那样）。\textbf{为何错}：混淆了“用扩散模型（model-free，要训练）”和“用扩散的结构/思想（model-based，免训练）”。MBD/DIAL-MPC 借的是扩散“多步去噪/退火”的\emph{结构}和“score 引导”的思想，而 score 由已知模型直接算，\emph{不训练任何网络、不需要任何数据}。\textbf{正解}：把它读作“用扩散的组织方式重新跑 MPPI”——零件是你已有的加权采样更新，新增的只是“多级退火”的调度和“用模型算 score”的视角。需要训练网络的那类（把扩散策略当 learned prior）是另一条前沿支线（见章末\nameref{sec:mv-frontier}）。
\end{pitfall}

\begin{practice}
\begin{enumerate}
  \item 对一个二次代价 $J=\tfrac12 U^\top H U$，实现 MPPI 单步加权更新，计算其更新方向与解析 score 上升方向 $-HU$ 的余弦相似度，验证 $\approx1$；再增大温度 $\lambda$ 或减小样本数 $K$，观察相似度如何变差（呼应 \cref{ch:mppi} 的 ESS）。
  \item 构造“宽浅局部阱 $+$ 深全局阱”的二维非凸地形，在\emph{相同样本预算}下对比单级 MPPI 与多级退火 MPPI 的平均最终代价与“命中全局阱”比例；再把退火调度方向写反（从小到大），验证退火失效。
  \item 用 \cref{eq:mv-tweedie} 与 \cref{eq:mv-mbd-posterior} 论证：MBD 的一级去噪，就是在“被噪声平滑过的目标分布”上做一步 MPPI。由此说明为什么 MBD 是“多级 MPPI”而非全新算法，以及为什么“多级降噪”是它逃离局部最优的关键（而“单步多跑几轮”不能）。
\end{enumerate}
\end{practice}

% ════════════════════════════════════════════════════════════════
\section{学习世界模型：TD-MPC 与 latent 空间规划}
\label{sec:mv-tdmpc}

\paragraph{动机：采样规划一直假设你“有一个能快速 rollout 的真实模型”。}
\cref{ch:mppi} 与本章前两节的规划循环都默认一个前提——\emph{你手上有一个能快速 rollout 的真实动力学模型}（刚体方程、MuJoCo 仿真器、或解析模型）。当你想把采样式 MPC 推得更远时，这个前提出现两道裂缝，而 \cref{sec:mv-diffusion} 的退火\emph{没有}弥合它们（退火仍在真实动力学里 rollout、仍是有限时域）：
\begin{itemize}
  \item \textbf{限制一：rollout 依赖真实动力学——它可能慢、贵，或根本不存在。} 高维全阶动力学、复杂接触求解、重型仿真器会压垮实时性；更极端的是观测为像素时\emph{根本没有可 rollout 的模型}（没有“像素的动力学方程”）。
  \item \textbf{限制二：短时域看不到长期——决策短视。} 规划时域 $H$ 受限（计算成本线性增长；用学习模型时 $H$ 越长累积误差越大）；若最优解需要“先忍受短期低回报、换取远期高回报”而那远期落在 $H$ 之外，短时域规划器就会短视。
\end{itemize}

\paragraph{核心观念：给采样规划换两个零件。}
\textbf{TD-MPC}（Hansen 等 $2022$）的答案听起来朴素、做起来精妙：\emph{不再依赖真实动力学，而从数据里学一个“世界模型”——既学一个轻量的 latent 动力学（补限制一），又学一个价值函数当终端价值（补限制二）}\cite{hansen2022tdmpc}。

\begin{table}[H]\centering\small
\caption{TD-MPC 相对 vanilla 采样 MPC：换了哪两个零件}
\label{tab:mv-twoparts}
\begin{tabular}{p{0.26\textwidth} p{0.30\textwidth} p{0.34\textwidth}}
\toprule
规划循环的环节 & vanilla MPPI/CEM（\cref{ch:mppi}） & TD-MPC（本节） \\
\midrule
在什么空间规划 & 真实状态空间 & \emph{学习的 latent 空间} \\
采样什么 & 动作序列（高斯） & 动作序列（高斯）——\emph{不变} \\
用什么 rollout & 真实动力学（仿真器/方程） & \emph{学习的 latent 动力学 $d_\theta$} \\
每步代价/reward & 真实代价函数 & \emph{学习的 reward 预测 $r_\theta$} \\
怎么更新分布 & 指数加权（MPPI）/精英（CEM） & 同左——\emph{不变} \\
时域末端怎么处理 & 直接截断（之后回报丢失） & \emph{加学习的终端价值 $Q_\theta$} \\
执行 & 第一个动作，滚动时域 & 同左——\emph{不变} \\
\bottomrule
\end{tabular}
\end{table}

\noindent 盯住“不变”那几行——\emph{采样什么、怎么更新、怎么执行：采样规划的内核一行没改}。变的只有两处：用学习的 latent 动力学（$+$ reward）替换真实动力学（补限制一：快、可批量、能吃像素），在末端加一个学习的价值（补限制二：短时域也有远见）。所以 TD-MPC 不是全新算法，而是把你已会的采样规划换上“学习的动力学”和“学习的终端价值”两个零件。

\subsection{三组件架构与 latent 空间规划}
\label{sec:mv-tdmpc-arch}

\begin{definition}[TD-MPC 世界模型的四个组件]\label{def:mv-worldmodel}
采样规划要在 latent 空间 rollout 一条动作序列、算回报、末端加价值，为此需要四个网络\cite{hansen2022tdmpc}：
\begin{itemize}
  \item \textbf{编码器} $z=h_\theta(o)$：把观测 $o$（低维状态向量或高维像素）压成紧凑 latent $z$；
  \item \textbf{latent 动力学} $z'=d_\theta(z,a)$：给定当前 latent 与动作，预测下一 latent，\emph{全程在 latent 空间、不碰原始观测}；
  \item \textbf{reward 预测} $\hat r=r_\theta(z,a)$：在 latent 上预测这一步 reward；
  \item \textbf{价值函数} $Q_\theta(z,a)$：估计“从 $z$ 做 $a$ 之后、一直到无穷远未来的期望折扣回报”。
\end{itemize}
前三样合称 TOLD（Task-Oriented Latent Dynamics，任务导向的 latent 动力学）。
\end{definition}

\begin{insight}
TOLD “任务导向”四字是 TD-MPC 区别于早期 latent world model（PlaNet/Dreamer）的命门。早期方法的 latent 是\emph{为重建观测}而学的——latent 要能解码回原始图像，必须编码图像里的一切（含无关细节如背景、光照）。TD-MPC 反其道而行：latent \emph{不为重建观测服务、只为预测 reward 和价值服务}，与回报无关的信息可完全丢弃。这是深刻的“少即是多”——\emph{做控制根本不需要一个能重建世界的模型，只需要一个能预测“与回报相关的未来”的模型}，于是不浪费一丝容量在无关重建上（这也是 TD-MPC 能用相对小的模型达到好性能的原因）。它呼应一切表示学习里“为下游任务定制表示、而非追求通用重建”的智慧——对比学习、BYOL 等自监督方法用的也是“在表示空间自洽、不重建输入”的思想。
\end{insight}

\noindent 为什么 latent 动力学用简单 MLP 而非 RNN/Transformer？其一，\emph{规划要求极快、极规整}——批量 rollout 成千上万条轨迹，MLP（纯矩阵乘 $+$ 激活）GPU 一次推完所有轨迹，RNN 的时序展开、Transformer 的注意力都会拖慢；其二，\emph{latent 已吸收“状态”信息}，动力学只需马尔可夫式单步映射，不需要 RNN 的记忆；其三，\emph{简单模型更好训、更稳、更省样本}。这也是为什么 TD-MPC 能直接从摄像头学控制而经典采样 MPC 不能——它把规划从“无法建模的像素空间”搬到“学出来的、可建模的 latent 空间”（像素只在“编码进 latent”那一步出现一次）。

\begin{algo}[TD-MPC 在 latent 空间的采样规划]\label{alg:mv-tdmpc-plan}
把 \cref{ch:mppi} 的 MPPI/CEM 搬进 latent 空间，规划循环几乎一字不变，只是 rollout 换成 latent 动力学、末端加价值：
\end{algo}

\begin{codebox}[Python]{TD-MPC latent 空间规划（核心循环）}
def plan(obs, prev_mean, model, H=5, K=512, iters=6, num_elites=64,
         num_pi_trajs=24, gamma=0.99, temperature=0.5):
    z = model.encode(obs)                    # 1. 编码: 观测 -> 初始 latent
    mean = shift(prev_mean)                  # 2. warm-start: 上周期解平移一格
    std  = torch.ones(H, action_dim) * 0.5
    for i in range(iters):                   # 3. 采样优化迭代(CEM/MPPI 式)
        noise   = torch.randn(K - num_pi_trajs, H, action_dim) * std
        sampled = (mean.unsqueeze(0) + noise).clamp(-1, 1)
        pi_acts = rollout_policy_prior(model, z, H)   # 一部分来自 policy prior
        actions = torch.cat([sampled, pi_acts], dim=0)
        z_k = z.unsqueeze(0).expand(K, -1)   # (K, latent_dim) 初始 latent 复制 K 份
        G, discount = torch.zeros(K), 1.0
        for t in range(H):                   # 在 latent 空间并行 rollout, 累积折扣 reward
            G   += discount * model.reward_(z_k, actions[:, t]).squeeze(-1)
            z_k  = model.next(z_k, actions[:, t])
            discount *= gamma
        a_H = model.pi_action(z_k)           # 终端价值(本节核心): 估 H 之后的长期回报
        G  += discount * model.Q(z_k, a_H).squeeze(-1)
        elite = G.topk(num_elites).indices   # 取精英(CEM)或指数加权(MPPI)更新分布
        w     = torch.softmax(G[elite] / temperature, dim=0)
        mean  = (w.view(-1,1,1) * actions[elite]).sum(0)
        std   = torch.sqrt((w.view(-1,1,1) *
                  (actions[elite] - mean.unsqueeze(0))**2).sum(0)).clamp(min=0.05)
    return mean[0]                           # 4. 执行第一个动作(滚动时域)
\end{codebox}

\noindent 与 \cref{ch:mppi} 的 vanilla MPPI 并排，结构\emph{完全同构}（采样—rollout—算回报—更新分布—迭代—执行第一个），差异只有三处：rollout 在 latent 空间（用 $d_\theta/r_\theta$）、末端加价值（$Q_\theta$）、采样掺入 \textbf{policy prior}。后者是一个与价值一起训练的策略网络 $\pi_\theta(z)$，提供“好动作大概在哪”的\emph{有信息}起点（纯高斯采样是盲目的、要靠大量采样 $+$ 迭代才找到好区域），让规划从好起点出发、快速收敛。这与 warm-start 一起，让 $6$ 次迭代足以收敛——呼应章末\nameref{sec:mv-frontier}提到的 learned prior 思想：用学习的先验引导采样、用在线优化保证质量。

\subsection{终端价值：短时域 + 学习的长期价值}
\label{sec:mv-tdmpc-value}

\paragraph{动机：本节最该带走的一个洞察。}
\cref{def:mv-worldmodel} 的前三个组件让你能在 latent 空间 rollout $H$ 步、算出这 $H$ 步的回报，但限制二说 $H$ 之后的长期回报看不到。直觉修法是“加长 $H$”，但加长有两个代价：计算线性增长，以及（用学习模型时）rollout 越长累积误差越大。TD-MPC 的答案完全不同：\emph{不加长 $H$，而在短 $H$ 的末端加一个学习的价值，把“$H$ 之外的长期回报”一次性估进来}。

\begin{theorem}[终端价值把无穷远未来塞进短时域]\label{thm:mv-terminal}
规划器 rollout $H$ 步、算出折扣回报，在末端 latent $z_H$ 处加上终端价值\cite{hansen2022tdmpc}：
\begin{equation}
  G=\underbrace{\sum_{t=0}^{H-1}\gamma^t\,r_\theta(z_t,a_t)}_{\text{$H$ 步内回报（rollout 算出）}}+\underbrace{\gamma^H\,Q_\theta(z_H,a_H)}_{\text{$H$ 之后长期回报（价值估出）}}.
  \label{eq:mv-terminal}
\end{equation}
第二项用学习的价值 $Q_\theta(z_H,a_H)$ 估计“从 $z_H$ 出发、之后一直走到无穷远的期望折扣回报”，把“$H$ 之后的所有回报”压成一个数。于是规划器虽只显式 rollout $H$ 步，却通过这个终端价值“看到”了 $H$ 之后的长期后果。注意前 $H$ 步用 reward 预测（\emph{零售}——一步一步算即时回报），末端换价值（\emph{批发}——把剩下无穷多步打包估成一个数）；若末端继续用 reward 只能再算“一步”，覆盖不了长期。终端价值\emph{必须}乘 $\gamma^H$ 才能和前面折扣一致。
\end{theorem}

\begin{insight}
终端价值是 TD-MPC 把“有限时域规划”和“无限时域最优”缝合的针脚。纯短时域规划求的是“未来 $H$ 步内的局部最优”，它在数学上\emph{不等于}真正的“无限时域全局最优”——除非把 $H$ 之后的所有回报也算进来，而这正是价值函数的定义。所以加上终端价值后，规划目标从“$H$ 步局部最优”变成对“无限时域最优”的一个\emph{有原则的近似}：用“短时域的计算成本”逼近“无限时域的决策质量”，前提是有一个够准的价值函数。这也解释了为什么 TD-MPC 必须\emph{同时}学好动力学和价值——动力学保证短 rollout 准、价值保证长期估计准，任一垮掉缝合就失效。这一招正是 TD-MPC（乃至 AlphaZero 的 MCTS $+$ value network）的命脉：\emph{有限前瞻 $+$ 学习价值}——用前瞻精确推演够得着的近期、用学习的价值概括够不着的远期。
\end{insight}

\noindent 为什么不\emph{纯靠}价值（像 SAC 那样直接用 $Q_\theta(z_0,a_0)$ 选动作）？因为价值是学出来的、有近似误差，而 rollout 那 $H$ 步是用模型实打实推演的、相对可靠。纯靠价值把全部信任押在一个有误差的函数上；TD-MPC 用“短 rollout 推演眼前 $+$ 价值补长期”分散风险——rollout 步数 $H$ 就是“信任模型推演 vs 信任价值”的旋钮，TD-MPC 选一个小的 $H$（$3\sim5$）取甜点。这也让“短 $H$”成为可能、间接绕开学习模型的累积误差：latent 动力学只 rollout $3\sim5$ 步，误差还没长大就停下，再交给价值。三个设计像三根互相支撑的柱子——latent 模型提供速度、短 $H$ 提供准确、终端价值提供远见，抽掉任何一根整座结构就塌。

\subsection{联合训练与 TD-MPC2 规模化}
\label{sec:mv-tdmpc-train}

\paragraph{联合训练：四网络一起学。}
TD-MPC 是一个完整的强化学习算法——用当前模型规划、与环境交互、把 transition $(o,a,r,o')$ 存进 replay buffer（\emph{规划即采集}）、从 buffer 采数据训练、模型变好后规划更好，循环往复（off-policy，数据可反复利用、样本高效）。核心是一个\emph{联合损失}，从 buffer 采一段长度 $H$ 的轨迹，同时训练四个网络\cite{hansen2022tdmpc}：
\begin{equation}
  \mathcal{L}=\sum_{t=0}^{H}\lambda^t\Big[\underbrace{c_1\|z_t-\mathrm{sg}(h_\theta(o_t))\|^2}_{\text{(1) latent 一致性}}+\underbrace{c_2(\hat r_t-r_t)^2}_{\text{(2) reward 预测}}+\underbrace{c_3(\hat Q_t-Q_{\text{target}})^2}_{\text{(3) TD / 价值}}\Big],
  \label{eq:mv-tdmpc-loss}
\end{equation}
其中（\textbf{此处 $\lambda$ 指时步衰减权重，非温度}）$z_t$ 由初始 latent 用动力学\emph{递推}得到（$z_{t+1}=d_\theta(z_t,a_t)$），$\mathrm{sg}$ 是 stop-gradient，TD 目标 $Q_{\text{target}}=r_t+\gamma Q_{\bar\theta}(h_\theta(o_{t+1}),a_{t+1}')$（目标网络 $Q_{\bar\theta}$ 参数缓慢跟随主网络，动作从 policy prior 采）。

\begin{insight}
(1) \textbf{latent 一致性损失}是 TD-MPC 最有特色、也是它区别于 Dreamer/PlaNet 的命门。它要求“用动力学\emph{递推}出来的 latent $z_t$，接近真实第 $t$ 步观测\emph{直接编码}出来的 latent $h_\theta(o_t)$”——\emph{在 latent 空间内部比较，根本不碰原始观测}。Dreamer/PlaNet 用\emph{重建损失}（要求 latent 能解码回原始观测），逼模型编码观测里的一切（含无关细节）；latent 一致性只要求“latent 层面动态预测一致”，latent 里编码什么由下游 reward/价值任务\emph{隐式决定}（任务导向），而非被“重建观测”这个与决策无关的目标绑架。注意右边的 stop-gradient：编码目标被当成固定靶子，梯度只更新“想象”一侧，防止编码器把所有 latent 塌到一个点（平凡解）。这是现代 model-based RL 从“重建式”走向“任务导向式”的技术核心，也是 TD-MPC 不需要解码器、模型能小而有效的原因。
\end{insight}

\noindent 价值用 \textbf{TD（时序差分）} bootstrap 学习——“当前 $Q$ 应等于这一步真实 reward $+$ 折扣 $\times$ 下一状态价值估计”，无需走到 episode 结束就能学到长期价值（对比蒙特卡洛“跑完整 episode”：TD 不必等结束、方差也小，代价是有偏，但偏差随训练逐步消除）。价值过估计（$Q$ 网络近似误差被 $\max$ 操作系统性放大——$\max$ 专挑被高估最多的动作）用 SAC 迁移来的\emph{双 $Q$/集成、取较小值}抑制。三项损失里 latent 一致性权重 $c_1$ 通常最大（它是地基，动力学预测不准则全盘皆错），价值权重 $c_3$ 通常最小（价值靠 bootstrap 学、目标本身在变，给小权重让它“慢慢学”）。

\paragraph{TD-MPC2：从单任务原型到通用控制器。}
TD-MPC（$2022$）是\emph{单任务}的、且每个任务往往要调超参。\textbf{TD-MPC2}（Hansen 等 $2024$）把它\emph{规模化、鲁棒化}：用单一超参集覆盖所有任务、并随模型/数据规模持续提升，在 $104$ 个在线 RL 任务、$4$ 个任务域上用单一超参集显著超越 SAC、DreamerV3 与原版 TD-MPC 等基线，并训了一个 $317$M 参数的 agent 执行 $80$ 个跨域、跨实体、跨动作空间的任务\cite{hansen2024tdmpc2}。关键改进：

\begin{table}[H]\centering\small
\caption{TD-MPC 到 TD-MPC2 的规模化改进}
\label{tab:mv-tdmpc2}
\begin{tabular}{p{0.24\textwidth} p{0.30\textwidth} p{0.36\textwidth}}
\toprule
改进 & 是什么 & 为何利于规模化 \\
\midrule
implicit/decoder-free & 彻底无解码器，连 latent 重建都不做 & 一个模型吃多种观测（状态/像素/异维）；“任务导向”做到极致 \\
SimNorm & latent 分块、每块做 softmax 归一化 & 把 latent 限制在概率单纯形上、范数有界，根除多步 rollout 的尺度发散 \\
离散回归 & reward/价值在对数空间分 bin、two-hot 编码、预测分布再解码 & 分类损失对目标量级不敏感，跨“reward 差几个数量级”的任务稳定 \\
$Q$ 集成 & \texttt{num\_q} 个 $Q$ 网络 $+$ dropout，随机取两个取较小值 & 任务越多过估计风险越大，集成更强地抑制 \\
任务嵌入 $+$ 动作 masking & 可学习任务向量拼到 latent；统一最大动作维、mask 无效维 & 一个网络处理不同观测/动作空间的任务（共性共享、特异性区分） \\
\bottomrule
\end{tabular}
\end{table}

\begin{insight}
TD-MPC2 \emph{没有改核心思想}（latent 模型 $+$ 终端价值 $+$ TD 还是那一套，\cref{def:mv-worldmodel,thm:mv-terminal,eq:mv-tdmpc-loss} 全部适用），它改的全是“让核心思想能规模化、鲁棒化”的工程。这揭示一个深刻规律：\emph{一个方法要真正可用、可规模化，光有好的核心思想不够，还需要一整套“稳定化、统一化”的工程改进}——SimNorm 稳 latent 尺度、离散回归稳 reward 量级、任务嵌入 $+$ masking 统一不同任务、$Q$ 集成稳价值。这些单看都是“工程细节”，合起来却是“能不能规模化”的分水岭。这个转变的意义，不亚于深度学习史上从“逐任务精调的小网络”到“一套 Transformer 架构 $+$ 大数据通吃”的转变——而 TD-MPC2 单任务默认模型才 $5$M 参数（任务导向、不浪费容量），需要大模型的是\emph{多任务}场景（一个模型装下几十上百个任务的知识）。
\end{insight}

\paragraph{TD-MPC 与 Dreamer 怎么选。}
学习世界模型这条线上最常拿来和 TD-MPC 比较的是 \textbf{Dreamer 系列}\cite{hafner2023dreamerv3}，两者在两个关键设计上分道扬镳：\emph{重建式 vs 任务导向}（Dreamer 的 latent 含解码器、能重建观测；TD-MPC 任务导向、decoder-free、只为预测回报）；\emph{想象优化 vs 在线规划}（Dreamer 用世界模型在“想象”里\emph{训练一个策略}、决策时直接用策略；TD-MPC 用世界模型做\emph{在线采样规划}，每步用 MPPI/CEM 在 latent 里搜索）。没有绝对赢家——Dreamer 决策快（训好策略后一次前向），TD-MPC 在线规划更灵活、短时域对模型误差更鲁棒。两者的设计哲学不同，各有适用场景。

\begin{pitfall}[编程陷阱：latent rollout 没向量化，用 Python 循环逐条 rollout K 条]
用一个 \texttt{for k in range(K)} 外层循环逐条 rollout。\textbf{为何错}：彻底浪费 latent 动力学“可大批量并行”的核心优势，规划慢几个数量级、实时性崩——latent 规划“快”的全部意义荡然无存。\textbf{正解}：把初始 latent 复制 $K$ 份成 batch（\texttt{z.\allowbreak expand(K, -1)}），rollout 循环只沿时域维度（$H$ 步）走，每步对整个 batch 做一次 \texttt{next}/\allowbreak\texttt{reward\_}；横向的 $K$ 条永远批量并行，只有纵向的 $H$ 步串行（正是 \cref{alg:mv-tdmpc-plan} 的写法）。
\end{pitfall}

\begin{practice}
\begin{enumerate}
  \item 用自己的话写出 vanilla 采样式 MPC 的两个限制，各举一个会触发它的机器人任务；说明 TD-MPC 的“学习的 latent 动力学”和“学习的终端价值”分别补哪个限制，以及为什么限制二不能简单靠“加长 $H$”解决。
  \item 复现“近处小峰 vs 远处大峰”demo：一维点 $x'=x+a$、reward 在近处 $x=2$ 有小峰、远处 $x=10$ 有大峰，验证“无价值时 $H=3/5/15/30$ 都困在小峰、$H=3$ 加（用 value iteration 预算的）价值就到达大峰”，据此理解 \cref{thm:mv-terminal}。
  \item 论证 latent 一致性损失（\cref{eq:mv-tdmpc-loss} 第(1)项）为何用“递推 latent 接近编码 latent”而非“重建观测”，以及 stop-gradient 为何能防止 latent 塌缩。再说明它与对比学习/BYOL“在表示空间自洽、不重建输入”的共通之处。
\end{enumerate}
\end{practice}

% ════════════════════════════════════════════════════════════════
\section*{本章前沿与开放问题}
\label{sec:mv-frontier}

三条主线仍在快速延伸，按方向分组：

\textbf{安全方向（向硬约束靠拢）。}\;\textbf{Shield-MPPI}（Yin 等 $2023$）在 rollout 中嵌入离散控制屏障函数（DCBF）安全惩罚与“局部修复”，把明显违反安全的样本挡掉，用远少于 MPPI 的样本（$<0.5\%$）在普通 CPU 上实时运行、向“硬安全”靠拢\cite{yin2023shield}；\textbf{DRA-MPPI}（Trevisan 等 $2025$）用蒙特卡洛实时估计数百条采样轨迹与多个动态障碍的\emph{联合碰撞概率}、缓解“机器人冻结问题”，专攻拥挤动态环境导航。

\textbf{规模化与统一框架。}\;承本章末尾——AIS（自适应重要性采样）、\textbf{Biased-MPPI}（Trevisan \& Alonso-Mora $2024$）把提议分布做成可在线自适应、可融合多个 ancillary 控制器作先验的更一般形式\cite{trevisan2024biased}；这类工作把“改提议的某一个侧面”推向“最优地构造提议分布”的统一框架（CEM 迭代精炼、CoVO 最优协方差、Gibbs 测度统一）。\textbf{Transformer 化}是另一条规模化伏笔：正如深度学习从“逐任务精调的小网络”走向“一套 Transformer 架构 $+$ 大数据通吃”（\cref{sec:mv-tdmpc-train} 的 TD-MPC2 单一超参 $+$ $317$M 模型正是 RL 版的这一步），把采样规划的提议/世界模型用 Transformer 规模化、用大规模离线数据训练，是把这套方法推向“通用控制器/具身大模型”的活跃方向。\pz{存疑：“Transformer-MPPI”作为一个独立命名方法在源讲义中未出现，此处按“规模化/Transformer 化趋势”归纳，具体代表工作待联网核实是否有专名。}

\textbf{learned prior 混合路线。}\;把扩散模型当 MPPI 的 learned prior 或近似动力学：\textbf{D-MPC}（Zhou 等 $2024$，DeepMind）用扩散学“多步动力学 $+$ 多步动作提议”，再用“采样—打分—排序”（sample-score-rank，即随机射击式）在线规划，可在运行时优化新奖励、适配新动力学\cite{zhou2024dmpc}；\textbf{Diffusion-MPC}（Huang 等 $2025$）把扩散当近似动力学先验、在去噪每步融入奖励规划与约束投影，直接作用于演化中的状态序列\cite{huang2025diffusionmpc}。回到 Gibbs 统一视角：MBD/DIAL 的 score 全由\emph{已知模型}算（纯 model-based 端），D-MPC/Diffusion-MPC 的 score 部分来自\emph{学习的先验}，TD-MPC 在中间偏数据一侧（学模型 $+$ 数据），它们只是在“score 从哪来”这条“模型$\leftrightarrow$数据”光谱上取了不同的点。

\textbf{核心开放问题：硬约束如何以数学严格的方式嵌入。}\;贯穿本章的一句话是“MPPI 的约束是软的”：RMPPI 给自由能增长上界、C²U-MPPI 给碰撞概率约束、Shield-MPPI 用 CBF 滤波——都在\emph{逼近}安全，但\emph{没有一个}能像传统鲁棒 MPC 那样在有界扰动下给出“硬约束必然满足”的严格保证。如何在保留 MPPI“处理非光滑代价与黑箱动力学”这一核心优势的同时获得硬约束严格满足，是采样式 MPC 与优化式 MPC 之间最后一道没完全打通的鸿沟。这些局限并非孤立，它们共根于——MPPI 是\emph{无导数、基于采样、单步闭式}的方法：正是“无导数”让它能处理黑箱和非光滑（优点）、也让它拿不到硬约束所需的梯度（局限）；要突破往往要“放弃一点 MPPI 的纯粹性”——引入迭代（CEM）、引入梯度（混合方法）、引入学习（TD-MPC/扩散先验）。

% ════════════════════════════════════════════════════════════════
\section*{本章常见误解汇总}
\begin{table}[H]\centering\small
\begin{tabular}{p{0.46\textwidth} p{0.46\textwidth}}
\toprule
\textbf{常见误解} & \textbf{正确理解} \\
\midrule
“鲁棒”就是把协方差调大、多探索 & 是\emph{中心（位置）}在闭环下失效，不是宽度不够（\cref{sec:mv-robust} 陷阱）\\
RMPPI $=$ Tube-MPPI $+$ 执行端反馈 & 必须把反馈写进 rollout（\cref{eq:mv-rmppi-rollout}），否则退化成换皮 Tube \\
自由能上界 $=$ 硬约束安全保证 & 是概率/能量意义的鲁棒性，仍是软约束（\cref{thm:mv-freenergy}）\\
事后低通滤波与 Smooth-MPPI 等价 & 平滑\emph{位置}不同：前者优化之后（破坏最优性），后者优化之内（保最优，\cref{thm:mv-lift}）\\
重尾 $=$ 把高斯方差 $\sigma$ 调大 & 重尾是“尖峰 $+$ 厚尾”，宽高斯是“整体摊平”，探索行为不同 \\
SVG-MPPI 同时覆盖所有 mode & 它是 mode-seeking（只锁一个目标 mode），非 mode-covering（\cref{alg:mv-svg}）\\
sigma 点 $=$ 少撒几个随机样本 & 是\emph{确定性}求积点（沿协方差矩阵\emph{平方根}摆放），非随机抽样 \\
扩散 MPC 要训练一个扩散网络 & MBD/DIAL 是 diffusion-\emph{inspired}、用模型算 score、\emph{免训练免数据}（\cref{thm:mv-mbd}）\\
MPPI $=$ 单步去噪 $\Rightarrow$ 扩散没带来新东西 & 新增是“多步复合”（突破天花板）与“model-based score”两件（\cref{sec:mv-diffusion}）\\
加长 $H$ 与加终端价值等效 & 加长 $H$ 又慢又救不了短视（demo $H{=}30$ 仍困小峰）；终端价值一次性补长期（\cref{thm:mv-terminal}）\\
TD-MPC 是全新的、与 MPPI 无关的 RL & 是采样规划内核 $+$ 换两个零件（latent 动力学 $+$ 终端价值，\cref{tab:mv-twoparts}）\\
TD-MPC 的 latent 要能重建观测 & 任务导向、decoder-free，只为预测 reward/价值（\cref{def:mv-worldmodel}）\\
latent 维度/模型越大越好 & 任务导向“够用就好”，单任务 $5$M 参数即可（\cref{sec:mv-tdmpc-train}）\\
\bottomrule
\end{tabular}
\end{table}

% ════════════════════════════════════════════════════════════════
\section*{本章小结}

\begin{table}[H]\centering\small
\caption{符号表}
\begin{tabular}{lll}
\toprule
符号 & 含义 & 首次出现 \\
\midrule
$U,\ \varepsilon_k,\ w_k$ & 控制序列 / rollout 噪声 / MPPI 权重 & \cref{eq:mv-mppi-update} \\
$S_k\ (J_k),\ \lambda,\ \rho$ & 第 $k$ 条 rollout 代价 / 温度 / 基线 $\min_k S_k$ & \cref{eq:mv-mppi-update} \\
$\bar x,\ K_{\text{fb}},\ x_k$ & 名义状态 / 反馈增益 / 第 $k$ 条 rollout 真实状态 & \cref{eq:mv-rmppi-rollout} \\
$\delta u,\ \Sigma_\delta$ & 控制导数 / 导数空间采样协方差 & \cref{eq:mv-inputlift} \\
$\pi,\ \mathrm{LogN}$ & NLN 高斯分量权重 / 对数正态 & \cref{eq:mv-nln} \\
$\theta_i,\ k(\cdot,\cdot)$ & SVGD 粒子（候选控制序列）/ RBF 核 & \cref{eq:mv-svgd} \\
$\chi_i,\ w_i^m,w_i^c,\ P$ & sigma 点 / 均值与协方差权重 / 协方差 & \cref{eq:mv-sigma} \\
$p_\sigma=p*\mathcal{N}(0,\sigma^2 I)$ & 目标分布与高斯核卷积（平滑版） & \cref{eq:mv-smoothed-score} \\
$\bar\alpha_k,\ \tau_k,\ \hat\tau_0$ & 噪声调度 / 带噪轨迹 / 后验均值 & \cref{thm:mv-mbd} \\
$h_\theta,d_\theta,r_\theta,Q_\theta$ & TD-MPC 编码器/latent 动力学/reward/价值 & \cref{def:mv-worldmodel} \\
$z,\ G,\ \gamma$ & latent 状态 / rollout 累积回报 / 折扣 & \cref{alg:mv-tdmpc-plan} \\
$\mathrm{sg}(\cdot),\ Q_{\text{target}}$ & stop-gradient / TD 目标 & \cref{eq:mv-tdmpc-loss} \\
\bottomrule
\end{tabular}
\end{table}

\begin{table}[H]\centering\small
\caption{定理速查表}
\begin{tabular}{lll}
\toprule
定理 / 公式 & 一句话说明 & 对应 \\
\midrule
RMPPI rollout \cref{eq:mv-rmppi-rollout} & 反馈入采样、权重用名义代价 $\Rightarrow$ 闭环有效 & \cref{sec:mv-robust} \\
自由能增长上界 \cref{thm:mv-freenergy} & 鲁棒 $=$ 约束满足 $+$ 跟踪性能 $+$ 采样误差三旋钮 & \cref{sec:mv-robust} \\
input-lifting \cref{eq:mv-inputlift,thm:mv-lift} & 导数空间采样 $=$ 扩状态 $\Rightarrow$ 平滑且保最优 & \cref{sec:mv-smooth} \\
NLN 重尾 \cref{eq:mv-nln} & 高斯主体 $+$ 对称化对数正态尾 $\Rightarrow$ 保可行下限 & \cref{sec:mv-log} \\
SVGD 寻峰 \cref{eq:mv-svgd,alg:mv-svg} & 反向 KL $\Rightarrow$ 锁单峰，避开多模平均撞墙 & \cref{sec:mv-svg} \\
UT 传播 \cref{eq:mv-sigma,eq:mv-ut-recombine} & $2n{+}1$ 确定性点传均值$+$协方差 $\Rightarrow$ 估碰撞概率 & \cref{sec:mv-umppi} \\
同构 \cref{thm:mv-isomorphism} & MPPI 一步 $=$ 去噪一步 $=$ score 上升一步 & \cref{sec:mv-bridge} \\
采样方差 $=$ 平滑尺度 \cref{eq:mv-smoothed-score} & 加权平均 $=$ 平滑分布的 score；退火 $=$ 由粗到精 & \cref{sec:mv-bridge} \\
MBD \cref{eq:mv-mbd-posterior,eq:mv-tweedie} & 模型直接算 score（后验均值）$\Rightarrow$ 免训练免数据 & \cref{sec:mv-mbd} \\
DIAL 退火 \cref{alg:mv-dial} & 多级退火 MPPI $\Rightarrow$ 全阶力矩控制 & \cref{sec:mv-dial} \\
换两个零件 \cref{tab:mv-twoparts} & latent 动力学 $+$ 终端价值 $\Rightarrow$ 内核不变 & \cref{sec:mv-tdmpc} \\
终端价值 \cref{eq:mv-terminal,thm:mv-terminal} & 短 $H$ rollout $+$ $\gamma^H Q$ $\Rightarrow$ 短时域也有远见 & \cref{sec:mv-tdmpc-value} \\
联合损失 \cref{eq:mv-tdmpc-loss} & latent 一致性 $+$ reward $+$ TD $\Rightarrow$ 任务导向世界模型 & \cref{sec:mv-tdmpc-train} \\
\bottomrule
\end{tabular}
\end{table}

\begin{table}[H]\centering\small
\caption{知识点总表}
\begin{tabular}{clll}
\toprule
\# & 知识点 & 核心要点 & 难度 \\
\midrule
1 & 四缺陷同根 & 朴素高斯：开环/白噪声/薄尾单峰/只传均值 & \(\bigstar\bigstar\) \\
2 & Tube-MPPI 双层 & nominal MPPI $+$ ancillary（iLQG）；planner--tracker 冲突 & \(\bigstar\bigstar\bigstar\) \\
3 & RMPPI & 增广状态 $+$ 反馈入采样；自由能增长上界 & \(\bigstar\bigstar\bigstar\bigstar\) \\
4 & Smooth-MPPI & input-lifting $=$ 扩状态；保最优；t-轴抖动 & \(\bigstar\bigstar\bigstar\) \\
5 & Log-MPPI & NLN 重尾、对称化；保可行性下限 & \(\bigstar\bigstar\) \\
6 & SVG-MPPI & SVGD 寻峰、反向 KL；侦察一个 mode & \(\bigstar\bigstar\bigstar\) \\
7 & U-MPPI / CoVO & UT 传协方差估碰撞概率；最优协方差设计 & \(\bigstar\bigstar\bigstar\) \\
8 & 扩散同构 & MPPI 一步 $=$ 去噪 $=$ score 上升 & \(\bigstar\bigstar\bigstar\bigstar\) \\
9 & 采样方差 $=$ 平滑尺度 & 加权平均 $=$ $p_\sigma$ 的 score；退火由粗到精 & \(\bigstar\bigstar\bigstar\bigstar\) \\
10 & MBD & 模型算 score（Tweedie）；免训练免数据 & \(\bigstar\bigstar\bigstar\bigstar\) \\
11 & DIAL-MPC & 多级退火；免求导 $+$ 免陷局部；全阶力矩 & \(\bigstar\bigstar\bigstar\bigstar\) \\
12 & TD-MPC 换零件 & latent 动力学 $+$ 终端价值；内核不变 & \(\bigstar\bigstar\bigstar\) \\
13 & 任务导向 latent & 不重建观测、latent 一致性；少即是多 & \(\bigstar\bigstar\bigstar\bigstar\) \\
14 & 终端价值 & 短 $H$ rollout $+$ $\gamma^H Q$；有限前瞻 $+$ 学习价值 & \(\bigstar\bigstar\bigstar\bigstar\) \\
15 & TD-MPC2 规模化 & decoder-free $+$ SimNorm $+$ 离散回归 $+$ 单一超参 & \(\bigstar\bigstar\bigstar\bigstar\) \\
\bottomrule
\end{tabular}
\end{table}

% ════════════════════════════════════════════════════════════════
\section*{延伸阅读}
\begin{itemize}
  \item \textbf{鲁棒 MPPI}：Williams 等\cite{williams2018robust}（Tube-MPPI 原始论文）、Gandhi 等\cite{gandhi2021robust}（RMPPI 与自由能增长上界、MPPI 家族首个性能保证）。
  \item \textbf{平滑 / 重尾 / 多模 / 不确定}：Kim 等\cite{kim2022smooth}（Smooth-MPPI / input-lifting）、Mohamed 等\cite{mohamed2022logmppi}（Log-MPPI）、Honda 等\cite{honda2024svgmppi}（SVG-MPPI）与 Lambert 等\cite{lambert2020svmpc}（SV-MPC，多模对照）、Mohamed 等\cite{mohamed2025umppi}（U-MPPI）、Yi 等\cite{yi2024covo}（CoVO-MPC 最优协方差与收敛理论）。
  \item \textbf{扩散启发}：Ho 等\cite{ho2020ddpm}（DDPM）、Song--Ermon\cite{song2019score}（score-based 生成）、Pan 等\cite{pan2024mbd}（MBD，免训练扩散轨迹优化）、Xue 等\cite{xue2025dial}（DIAL-MPC，扩散退火全阶力矩控制）、Li 等\cite{li2025unified}（MPPI/RL/扩散的 Gibbs 统一）。
  \item \textbf{学习世界模型}：Hansen 等\cite{hansen2022tdmpc}（TD-MPC）、Hansen 等\cite{hansen2024tdmpc2}（TD-MPC2 规模化）、Hafner 等\cite{hafner2023dreamerv3}（DreamerV3，重建式对照）、Chi 等\cite{chi2023diffusion}（Diffusion Policy，learned prior 路线源头）。
  \item \textbf{前置工具}：Tassa 等\cite{tassa2012ilqg}（iLQG 稳定化，Tube/RMPPI 下层反馈来源）、Julier--Uhlmann\cite{julier1997unscented}（Unscented Transform，U-MPPI 的传播工具）。
\end{itemize}

\section*{本章与后续章节的关系}
\begin{table}[H]\centering\small
\begin{tabular}{lll}
\toprule
关联章节 & 关系 & 本章铺垫 \\
\midrule
\cref{ch:mppi} MPPI 基础 & 本章是其直接延续 & 加权更新 / 自由能 / ESS / warm-start \\
\cref{ch:planning} 运动规划导论 & 采样式 vs 优化式范式、构型空间 & 采样式 MPC 是局部优化器、需全局引导 \\
\cref{ch:lie} 李群与李代数 & 状态含姿态时的 rollout 插值/距离 & $\mathrm{Exp}/\mathrm{Log}$、测地线 \\
\bottomrule
\end{tabular}
\end{table}

\section*{故障排查手册}
\begin{enumerate}
  \item \textbf{症状}：仿真好、真车冲出安全边界。\textbf{排查}：是否“反馈只挂在执行端、没进 rollout”（退化成 Tube，\cref{sec:mv-robust} 陷阱）；自检——去掉 \cref{eq:mv-rmppi-rollout} 的反馈项后结果若无差别，说明反馈根本没进采样。\textbf{对策}：rollout 写成 $u_{\text{nom}}+K_{\text{fb}}(x_k-\bar x)+\varepsilon_k$、用真实动力学前向、权重用名义代价。
  \item \textbf{症状}：控制高频颤动、伤执行器。\textbf{排查}：是在控制空间撒白噪声（t-轴抖动）。\textbf{对策}：input-lifting 在导数空间采样再\emph{累积}积分 $u(t)=u(t{-}1)+\delta u(t)\,\mathrm{d}t$（漏 $\mathrm{d}t$ 或漏累积则低通失效）；勿用事后低通（破坏最优性 $+$ 相位滞后）。
  \item \textbf{症状}：杂乱环境 MPPI 输出近乎随机、卡死。\textbf{排查}：先查“有没有可行样本”（可行性，非效率）。\textbf{对策}：NLN 重尾捞回可行 rollout（确认对数正态分量已对称化）；必要时配 GP-MPPI 子目标引导方向。
  \item \textbf{症状}：双峰避障取了中间、撞墙。\textbf{原因}：单峰高斯加权平均落两峰间（\cref{sec:mv-svg}）。\textbf{对策}：SVGD 寻峰锁一个 mode、或轨迹聚类先分簇再各自加权；勿靠调 $\lambda$/$\Sigma$/加样本（表达力问题非参数问题）。
  \item \textbf{症状}：退火 MPC 效果不如单级。\textbf{排查}：温度/方差调度是否\emph{单调递减}（写反成由精到粗则失效）；是否只缩温度不缩方差（后期收不住）。\textbf{对策}：温度与方差都单调递减、高温级管全局、低温级管局部；起始 $\sigma_0$ 取到能覆盖问题全局尺度、$\mathrm{decay}$ 让末级方差降到精修所需小尺度。
  \item \textbf{症状}：latent 多步 rollout 数值发散、训练崩。\textbf{原因}：latent 尺度在递推中失控（\cref{sec:mv-tdmpc-train}）。\textbf{对策}：动力学输出过 SimNorm（把 latent 限制在概率单纯形上、范数有界）；reward/价值用离散回归稳量级。
  \item \textbf{症状}：TD-MPC 训练初期规划行为随机、不见好。\textbf{原因}：初期是随机动作\emph{预热}（seed steps）填 buffer $+$ 模型刚起步，行为随机是预期的、非 bug。\textbf{对策}：耐心等预热结束 $+$ 模型积累数据；若长期不收敛，查价值是否过估计（双 $Q$/集成）、latent 一致性损失是否降下去。
\end{enumerate}
